% =============================================================================
%  trpc-meeting.tex — เอกสารฉบับที่ 2 ของชุด "สัญญา tRPC" ของโครงการ ULMs
%
%  วาระประชุม 19 วาระ เรียงตามลำดับว่าอะไรบล็อกอะไร
%  ใช้เป็นเอกสารตั้งต้นในห้องประชุม ผู้เข้าร่วมทุกคนควรอ่านภาค 1 มาก่อน
%
%  เอกสารคู่กัน: trpc-guide.tex อธิบายว่าทำไมแต่ละวาระถึงสำคัญ
%
%  ต้องคอมไพล์ด้วย XeLaTeX เท่านั้น
%      cd docs && latexmk trpc-meeting.tex     → build/trpc-meeting.pdf
%
%  ที่มาของข้อเท็จจริง: เหมือนกับ trpc-guide.tex (ดูภาคผนวก ง ของฉบับนั้น)
% =============================================================================

% =============================================================================
%  trpc-preamble.tex — preamble ร่วมของเอกสารชุด "สัญญา tRPC" ของโครงการ ULMs
%
%  ใช้โดย:  trpc-guide.tex    (ฉบับที่ 1 — หลักการ tRPC + ความเสี่ยงใน repo นี้)
%           trpc-meeting.tex  (ฉบับที่ 2 — วาระประชุมเรียงตามลำดับความสำคัญ)
%
%  ทำไมไม่ใช้ ulms-preamble.tex ที่มีอยู่แล้ว
%  ----------------------------------------
%  ulms-preamble.tex ตั้งใจให้เป็น "โทนเดียว (monochrome)" สำหรับเอกสารที่ต้อง
%  ส่งอาจารย์/พิมพ์ขาวดำ แต่เอกสารสอนสองฉบับนี้ใช้ "สีเชิงความหมาย" เป็นเครื่องมือ
%  สอนหลัก (สีเดียวกัน = แนวคิดเดียวกัน ทั้งในเนื้อความ ในไดอะแกรม และในตาราง)
%  จึงต้องมี preamble ของตัวเอง — ไม่ใช่ก๊อปมาแก้สีทับ
%
%  เอกสารที่เรียกใช้ต้องทำสองอย่างหลัง % =============================================================================
%  trpc-preamble.tex — preamble ร่วมของเอกสารชุด "สัญญา tRPC" ของโครงการ ULMs
%
%  ใช้โดย:  trpc-guide.tex    (ฉบับที่ 1 — หลักการ tRPC + ความเสี่ยงใน repo นี้)
%           trpc-meeting.tex  (ฉบับที่ 2 — วาระประชุมเรียงตามลำดับความสำคัญ)
%
%  ทำไมไม่ใช้ ulms-preamble.tex ที่มีอยู่แล้ว
%  ----------------------------------------
%  ulms-preamble.tex ตั้งใจให้เป็น "โทนเดียว (monochrome)" สำหรับเอกสารที่ต้อง
%  ส่งอาจารย์/พิมพ์ขาวดำ แต่เอกสารสอนสองฉบับนี้ใช้ "สีเชิงความหมาย" เป็นเครื่องมือ
%  สอนหลัก (สีเดียวกัน = แนวคิดเดียวกัน ทั้งในเนื้อความ ในไดอะแกรม และในตาราง)
%  จึงต้องมี preamble ของตัวเอง — ไม่ใช่ก๊อปมาแก้สีทับ
%
%  เอกสารที่เรียกใช้ต้องทำสองอย่างหลัง % =============================================================================
%  trpc-preamble.tex — preamble ร่วมของเอกสารชุด "สัญญา tRPC" ของโครงการ ULMs
%
%  ใช้โดย:  trpc-guide.tex    (ฉบับที่ 1 — หลักการ tRPC + ความเสี่ยงใน repo นี้)
%           trpc-meeting.tex  (ฉบับที่ 2 — วาระประชุมเรียงตามลำดับความสำคัญ)
%
%  ทำไมไม่ใช้ ulms-preamble.tex ที่มีอยู่แล้ว
%  ----------------------------------------
%  ulms-preamble.tex ตั้งใจให้เป็น "โทนเดียว (monochrome)" สำหรับเอกสารที่ต้อง
%  ส่งอาจารย์/พิมพ์ขาวดำ แต่เอกสารสอนสองฉบับนี้ใช้ "สีเชิงความหมาย" เป็นเครื่องมือ
%  สอนหลัก (สีเดียวกัน = แนวคิดเดียวกัน ทั้งในเนื้อความ ในไดอะแกรม และในตาราง)
%  จึงต้องมี preamble ของตัวเอง — ไม่ใช่ก๊อปมาแก้สีทับ
%
%  เอกสารที่เรียกใช้ต้องทำสองอย่างหลัง \input{trpc-preamble}:
%    1) \hypersetup{pdftitle={...}}
%    2) \begin{document}
%
%  Build:  cd docs && latexmk trpc-guide.tex     (.latexmkrc บังคับ xelatex ให้แล้ว)
%          PDF ออกที่ docs/build/ แล้วค่อยคัดลอกไป docs/report/
% =============================================================================

\documentclass[11pt,a4paper]{article}

% ---------------------------------------------------------------------
%  FONTS
%
%  ต้องใช้ XeLaTeX เท่านั้น — pdflatex เรียงพิมพ์ภาษาไทยไม่ได้เลย (ไม่ใช่ "ได้แต่ไม่สวย")
%
%  ใช้ตระกูล TLWG ที่มีทั้งอักษรไทยและละตินอยู่ในไฟล์เดียว จงใจเลี่ยงการจับคู่
%  ฟอนต์ละตินกับฟอนต์ไทยล้วนผ่าน ucharclasses เพราะ ucharclasses สลับฟอนต์ตาม
%  "บล็อกยูนิโคด" ของตัวอักษร โดยไม่รู้จักขอบเขตของ TeX group → ข้อความละตินที่
%  ตามหลัง \textbf{...} จะยังค้างอยู่ในฟอนต์ไทยซึ่งไม่มีสัญลักษณ์ละติน แล้ว
%  "หายไปจาก PDF เงียบ ๆ" โดยไม่มี warning
% ---------------------------------------------------------------------
\usepackage{fontspec}

\setmainfont{Laksaman}[Ligatures=TeX]          % ไทย+ละติน — เนื้อความ
\setsansfont{Garuda}[Ligatures=TeX]            % ไทย+ละติน หนักกว่า — หัวข้อ/ป้ายในไดอะแกรม
\setmonofont{DejaVu Sans Mono}[Scale=0.84]     % 0.84 ทำให้ x-height ใกล้เคียง TLWG

% --- การตัดบรรทัดภาษาไทย ---
% ภาษาไทยเขียนติดกันไม่มีช่องว่าง TeX จึงหาจุดตัดบรรทัดเองไม่ได้ และจะปล่อยให้
% บรรทัดล้นออกนอกหน้ากระดาษ สองบรรทัดนี้ยกงานให้ ICU ตัดคำแบบพจนานุกรม
% ค่า stretch ต้องมากกว่าศูนย์อย่างมีนัยสำคัญ ไม่งั้น TeX แทบไม่มีอิสระในการจัด
% บรรทัดไทยและจะออกมาเป็น overfull box แทนที่จะยืด/หดช่องไฟ
\XeTeXlinebreaklocale "th"
\XeTeXlinebreakskip = 0pt plus 0.15em minus 0.02em
\emergencystretch=2.2em

% หัวข้อ ป้ายในไดอะแกรม หัวตาราง header/footer ใช้ sans — เนื้อความไม่ใช้
% นี่คือสิ่งที่แยก "โครงสร้างสำหรับกวาดสายตา" ออกจาก "เนื้อหาสำหรับอ่าน"
\newcommand{\sansall}{\sffamily}

% ---------------------------------------------------------------------
%  LAYOUT
% ---------------------------------------------------------------------
\usepackage[a4paper,margin=2.3cm,top=2.5cm,bottom=2.5cm,headsep=12pt]{geometry}
\usepackage{amsmath,amssymb}
\usepackage{graphicx}
\graphicspath{{figures/}}
\usepackage{booktabs}
\usepackage{array}
\usepackage{longtable}
\usepackage{multirow}
\usepackage{xcolor}
\usepackage{tikz}
\usetikzlibrary{arrows.meta,positioning,calc,shapes.geometric,shapes.misc,
                fit,backgrounds,decorations.pathmorphing}
\usepackage[most]{tcolorbox}
\usepackage[font=small,labelfont=bf,width=0.95\textwidth]{caption}
\usepackage{fancyhdr}
\usepackage{enumitem}
\usepackage{titlesec}
\usepackage{float}        % figure[H] — ตรึงรูปไว้ข้างเนื้อความที่อธิบายมัน
\usepackage{microtype}
\usepackage{listings}
\usepackage[colorlinks=true,allcolors=ctr,breaklinks=true,
            bookmarksnumbered=true]{hyperref}

% ภาษาไทยต้องการระยะบรรทัดมากกว่าละติน สระบนซ้อนวรรณยุกต์ และมีสระล่างด้วย
% ที่ระยะปกติจะชนกันข้ามบรรทัด — 1.28 คือค่าที่ทดสอบแล้ว
\linespread{1.28}
\setlength{\parskip}{0.35em}
\setlength{\parindent}{0pt}

% ---------------------------------------------------------------------
%  สี — ระบบสีเชิงความหมาย (semantic colour)
%
%  ฐานสีคือชุด Okabe–Ito ซึ่งปลอดภัยสำหรับผู้ที่มองสีบางคู่ไม่แยก (colour-blind safe)
%
%  สีสี่สีหลักถูก "ตั้งชื่อกลาง ๆ" ไว้ที่นี่ แล้วให้แต่ละเอกสารประกาศความหมายของ
%  ตัวเองในหน้า "เอกสารนี้อ่านอย่างไร":
%
%    trpc-guide.tex   → 4 ชั้นของระบบ:   สัญญา / backend / frontend / เครือข่าย
%    trpc-meeting.tex → 4 ระดับเร่งด่วน: P0 / P1 / P2 / P3
%
%  สองเอกสารใช้คนละความหมายได้ เพราะแต่ละฉบับ "ประกาศระบบสีของตัวเองไว้หน้าแรก"
%  และคงเส้นคงวาภายในเล่ม — สิ่งที่ห้ามคือเปลี่ยนความหมายกลางเล่ม
% ---------------------------------------------------------------------
\definecolor{ctr}{HTML}{0072B2}      % ฟ้า   — สัญญา/schema/type   · หรือ P2
\definecolor{beo}{HTML}{D55E00}      % ส้ม   — ฝั่ง backend         · หรือ P1
\definecolor{feg}{HTML}{009E73}      % เขียว — ฝั่ง frontend        · หรือ P3
\definecolor{ops}{HTML}{7B5EA7}      % ม่วง  — เครือข่าย/รันไทม์/ops

\definecolor{ink}{HTML}{1A2733}      % เข้ม — หัวข้อรอง
\definecolor{muted}{HTML}{5B6B7A}    % ข้อความรอง เส้นคั่น ป้ายกำกับ
\definecolor{rule}{HTML}{C9D2DB}     % เส้นบาง ขอบไดอะแกรม
\definecolor{wash}{HTML}{F1F6FA}     % พื้นกล่องฟ้า/กลาง
\definecolor{warnwash}{HTML}{FDF3E7} % พื้นกล่องส้ม
\definecolor{goodwash}{HTML}{E9F4F0} % พื้นกล่องเขียว
\definecolor{badwash}{HTML}{FBEDED}  % พื้นกล่องแดง
\definecolor{bad}{HTML}{B3261E}      % แดง — ผิด/อันตราย · หรือ P0
\definecolor{codebg}{HTML}{F7F9FB}   % พื้นบล็อกโค้ด

% ---------------------------------------------------------------------
%  HEADINGS
%  \part ใช้เลขอารบิก เพื่อให้อ้างอิงข้ามเล่มอ่านว่า "ภาค 3" ไม่ใช่ "ภาค III"
% ---------------------------------------------------------------------
\renewcommand{\thepart}{\arabic{part}}
\titleformat{\part}[display]
  {\sansall\huge\bfseries\color{ctr}}
  {\normalsize\color{muted}ภาค \thepart}{0.4em}{}
\titlespacing*{\part}{0pt}{0pt}{1.5em}

\titleformat{\section}
  {\sansall\Large\bfseries\color{ctr!88!black}}{\thesection}{0.7em}{}
\titleformat{\subsection}
  {\sansall\large\bfseries\color{ink}}{\thesubsection}{0.6em}{}
\titleformat{\subsubsection}
  {\sansall\normalsize\bfseries\color{muted}}{\thesubsubsection}{0.55em}{}
\titlespacing*{\section}{0pt}{1.9em}{0.7em}
\titlespacing*{\subsection}{0pt}{1.3em}{0.45em}
\titlespacing*{\subsubsection}{0pt}{1.0em}{0.35em}

% หัวกระดาษ: ชื่อเอกสารคงที่ทางซ้าย ชื่อหัวข้อปัจจุบันทางขวา ผู้อ่านจะรู้ตลอดว่าอยู่ตรงไหน
% \docshort ถูกกำหนดในไฟล์เอกสารแต่ละฉบับ
\providecommand{\docshort}{}
\pagestyle{fancy}
\fancyhf{}
\renewcommand{\headrulewidth}{0.4pt}
\fancyhead[L]{\footnotesize\sansall\color{muted}\docshort}
\fancyhead[R]{\footnotesize\sansall\color{muted}\leftmark}
\fancyfoot[C]{\footnotesize\color{muted}\thepage}
\renewcommand{\sectionmark}[1]{\markboth{#1}{}}

% รายการแบบกระชับ — ค่าเริ่มต้นของ LaTeX ดูหลวมเมื่ออยู่กับ \parskip
\setlist[itemize]{itemsep=1pt,topsep=3pt,parsep=0pt,leftmargin=1.4em}
\setlist[enumerate]{itemsep=1pt,topsep=3pt,parsep=0pt,leftmargin=1.6em}
\setlist[description]{itemsep=2pt,topsep=3pt,parsep=0pt}

% ---------------------------------------------------------------------
%  กล่องข้อความ 4 แบบ ความหมายตายตัว ห้ามคิดแบบที่ห้ากลางเล่ม
%
%  colbacktitle ต้องตั้งให้เท่ากับ colback มิฉะนั้น tcolorbox จะใช้ colframe
%  (สีเข้ม) เป็นพื้นหลังของหัวกล่อง แล้วตัวอักษรหัวกล่องสีเข้มจะอ่านไม่ออก
% ---------------------------------------------------------------------
\tcbset{guidebox/.style={
  enhanced, breakable, boxrule=0pt, leftrule=2.5pt, arc=1pt,
  left=8pt, right=8pt, top=6pt, bottom=6pt,
  fonttitle=\sansall\bfseries\small, titlerule=0pt,
  toptitle=1pt, bottomtitle=1pt}}

% ฟ้า — ใจความที่ต้องจำ / ประโยคที่ร้อยหลายบทเข้าด้วยกัน
\newtcolorbox{keybox}[1][]{guidebox, colback=wash, colframe=ctr,
  colbacktitle=wash, coltitle=ctr!85!black, #1}
% ส้ม — กับดัก หรือข้อสมมติที่ถ้าลืมแล้วข้อสรุปจะผิด
\newtcolorbox{warnbox}[1][]{guidebox, colback=warnwash, colframe=beo,
  colbacktitle=warnwash, coltitle=beo!80!black, #1}
% เขียว — สิ่งที่ตรวจสอบกับของจริงในรีโปแล้ว
\newtcolorbox{goodbox}[1][]{guidebox, colback=goodwash, colframe=feg,
  colbacktitle=goodwash, coltitle=feg!65!black, #1}
% แดง — วิธีที่ผิดจริง ๆ อย่าทำ
\newtcolorbox{badbox}[1][]{guidebox, colback=badwash, colframe=bad,
  colbacktitle=badwash, coltitle=bad, #1}

% ---------------------------------------------------------------------
%  SEMANTIC MARKUP
% ---------------------------------------------------------------------
% ศัพท์อังกฤษที่แทรกในประโยคภาษาไทย ฟอนต์เนื้อความครอบคลุมละตินอยู่แล้วจึงไม่ต้อง
% สลับฟอนต์ — เก็บแมโครไว้เป็น "เครื่องหมายเชิงความหมาย" เผื่อวันหลังอยากเปลี่ยน
% การแสดงผลทั้งเล่มพร้อมกัน
\newcommand{\en}[1]{#1}
% ชื่อไฟล์ ฟังก์ชัน หรือค่าที่ปรากฏในโค้ดจริง
% ไม่บังคับขนาดตัวอักษร ปล่อยให้รับขนาดจากบริบท (สำคัญในตารางที่ใช้ \footnotesize
% มิฉะนั้นโค้ดจะกลายเป็นตัวใหญ่กว่าข้อความรอบข้างแล้วดันคอลัมน์ล้น)
% ฟอนต์ mono ถูกย่อไว้ 0.84 แล้วใน \setmonofont จึงพอดีกับเนื้อความอยู่แล้ว
\newcommand{\code}[1]{{\ttfamily #1}}

% ชื่อไฟล์ที่มีอักษรไทยอยู่ในชื่อ — ห้ามใช้ \code เพราะ DejaVu Sans Mono ไม่มี
% สัญลักษณ์ไทย ตัวอักษรจะหายไปจาก PDF โดยขึ้นแค่ warning "Missing character"
% ใช้ Garuda (sans ที่มีทั้งไทยและละติน) แทน
\newcommand{\fnth}[1]{{\sffamily\small #1}}

% ป้ายชื่อแนวคิด — สีเดียวกับที่ใช้ในไดอะแกรมและตาราง
\newcommand{\stg}[2]{\textcolor{#1}{\textbf{\en{#2}}}}

% ---------------------------------------------------------------------
%  ไวยากรณ์ของไดอะแกรม — นิยามครั้งเดียว ทุกรูปจะดูเป็นชุดเดียวกัน
%  หมายเหตุ: node ที่มีข้อความไทยต้องเผื่อ inner ysep มากกว่าไดอะแกรมภาษาละติน
%  เพราะสระบน/ล่างกินที่แนวตั้ง
% ---------------------------------------------------------------------
\tikzset{
  blk/.style={draw=rule, line width=0.6pt, rounded corners=2.5pt,
              fill=wash, minimum height=9mm, inner xsep=6pt, inner ysep=5pt,
              align=center, font=\sansall\footnotesize},
  stage/.style={blk, fill=ctr!8, draw=ctr!45},
  bestage/.style={blk, fill=beo!8, draw=beo!45},
  festage/.style={blk, fill=feg!8, draw=feg!45},
  opstage/.style={blk, fill=ops!8, draw=ops!45},
  ar/.style={-{Stealth[length=2.2mm,width=1.8mm]}, line width=0.7pt,
             draw=muted},
  lbl/.style={font=\sansall\scriptsize, text=muted, align=center},
  tag/.style={font=\sansall\scriptsize\bfseries, align=center},
}

% ---------------------------------------------------------------------
%  ตาราง
%  คอลัมน์แบบชิดซ้าย: ข้อความจัดชิดขอบสองข้าง (justified) ในคอลัมน์แคบจะเปิด
%  ช่องไฟกว้างมาก และแย่เป็นพิเศษกับภาษาไทย เพราะ ICU ตัดไทยเป็นก้อนที่ใหญ่กว่า
%  คำภาษาอังกฤษ → ใช้ L{} แทน p{} ในทุกคอลัมน์ที่มีข้อความยาว
% ---------------------------------------------------------------------
\newcolumntype{L}[1]{>{\raggedright\arraybackslash}p{#1}}
\captionsetup{skip=6pt}

% ---------------------------------------------------------------------
%  บล็อกโค้ด
%
%  ข้อจำกัดที่ต้องรู้: listings อ่านไฟล์แบบไบต์ ไม่เข้าใจ UTF-8 หลายไบต์
%  → "ห้ามใส่ภาษาไทยในบล็อกโค้ด" คอมเมนต์ในโค้ดทุกอันเป็นภาษาอังกฤษ
%    คำอธิบายภาษาไทยให้อยู่นอกบล็อกเสมอ
% ---------------------------------------------------------------------
\lstdefinelanguage{TypeScript}{
  keywords={await, async, break, case, catch, class, const, continue, default,
            delete, do, else, enum, export, extends, false, finally, for, from,
            function, if, implements, import, in, instanceof, interface, let,
            new, null, of, private, public, readonly, return, super, switch,
            this, throw, true, try, type, typeof, undefined, var, void, while,
            yield, satisfies, as},
  keywordstyle=\color{ctr}\bfseries,
  ndkeywords={string, number, boolean, Date, Promise, z, Record, Array},
  ndkeywordstyle=\color{ops},
  sensitive=true,
  comment=[l]{//},
  morecomment=[s]{/*}{*/},
  commentstyle=\color{muted}\itshape,
  stringstyle=\color{feg!70!black},
  morestring=[b]',
  morestring=[b]",
  morestring=[b]`,
}

\lstset{
  language=TypeScript,
  basicstyle=\ttfamily\scriptsize,
  backgroundcolor=\color{codebg},
  frame=leftline,
  framerule=2pt,
  rulecolor=\color{rule},
  framesep=7pt,
  xleftmargin=10pt,
  breaklines=true,
  breakatwhitespace=true,
  columns=fullflexible,
  keepspaces=true,
  showstringspaces=false,
  aboveskip=0.9em,
  belowskip=0.9em,
  literate={->}{{$\rightarrow$}}2 {=>}{{=>}}2,
}

% ชื่อ float ภาษาไทย — ถ้าไม่ตั้ง จะได้ "Figure 3" อยู่ใต้เนื้อความภาษาไทย
\renewcommand{\figurename}{รูปที่}
\renewcommand{\tablename}{ตารางที่}
\renewcommand{\refname}{เอกสารอ้างอิง}
\renewcommand{\contentsname}{สารบัญ}
% หมายเหตุ: ห้ามใส่ \color ใน \contentsname เพราะ \tableofcontents ส่งมันผ่าน
% \MakeUppercase ซึ่งจะเปลี่ยน ctr เป็น CTR แล้ว xcolor จะฟ้องว่าไม่รู้จักสีนั้น
% ถ้าจะจัดสไตล์หัวสารบัญ ให้ทำผ่าน \titleformat{\section} แทน
:
%    1) \hypersetup{pdftitle={...}}
%    2) \begin{document}
%
%  Build:  cd docs && latexmk trpc-guide.tex     (.latexmkrc บังคับ xelatex ให้แล้ว)
%          PDF ออกที่ docs/build/ แล้วค่อยคัดลอกไป docs/report/
% =============================================================================

\documentclass[11pt,a4paper]{article}

% ---------------------------------------------------------------------
%  FONTS
%
%  ต้องใช้ XeLaTeX เท่านั้น — pdflatex เรียงพิมพ์ภาษาไทยไม่ได้เลย (ไม่ใช่ "ได้แต่ไม่สวย")
%
%  ใช้ตระกูล TLWG ที่มีทั้งอักษรไทยและละตินอยู่ในไฟล์เดียว จงใจเลี่ยงการจับคู่
%  ฟอนต์ละตินกับฟอนต์ไทยล้วนผ่าน ucharclasses เพราะ ucharclasses สลับฟอนต์ตาม
%  "บล็อกยูนิโคด" ของตัวอักษร โดยไม่รู้จักขอบเขตของ TeX group → ข้อความละตินที่
%  ตามหลัง \textbf{...} จะยังค้างอยู่ในฟอนต์ไทยซึ่งไม่มีสัญลักษณ์ละติน แล้ว
%  "หายไปจาก PDF เงียบ ๆ" โดยไม่มี warning
% ---------------------------------------------------------------------
\usepackage{fontspec}

\setmainfont{Laksaman}[Ligatures=TeX]          % ไทย+ละติน — เนื้อความ
\setsansfont{Garuda}[Ligatures=TeX]            % ไทย+ละติน หนักกว่า — หัวข้อ/ป้ายในไดอะแกรม
\setmonofont{DejaVu Sans Mono}[Scale=0.84]     % 0.84 ทำให้ x-height ใกล้เคียง TLWG

% --- การตัดบรรทัดภาษาไทย ---
% ภาษาไทยเขียนติดกันไม่มีช่องว่าง TeX จึงหาจุดตัดบรรทัดเองไม่ได้ และจะปล่อยให้
% บรรทัดล้นออกนอกหน้ากระดาษ สองบรรทัดนี้ยกงานให้ ICU ตัดคำแบบพจนานุกรม
% ค่า stretch ต้องมากกว่าศูนย์อย่างมีนัยสำคัญ ไม่งั้น TeX แทบไม่มีอิสระในการจัด
% บรรทัดไทยและจะออกมาเป็น overfull box แทนที่จะยืด/หดช่องไฟ
\XeTeXlinebreaklocale "th"
\XeTeXlinebreakskip = 0pt plus 0.15em minus 0.02em
\emergencystretch=2.2em

% หัวข้อ ป้ายในไดอะแกรม หัวตาราง header/footer ใช้ sans — เนื้อความไม่ใช้
% นี่คือสิ่งที่แยก "โครงสร้างสำหรับกวาดสายตา" ออกจาก "เนื้อหาสำหรับอ่าน"
\newcommand{\sansall}{\sffamily}

% ---------------------------------------------------------------------
%  LAYOUT
% ---------------------------------------------------------------------
\usepackage[a4paper,margin=2.3cm,top=2.5cm,bottom=2.5cm,headsep=12pt]{geometry}
\usepackage{amsmath,amssymb}
\usepackage{graphicx}
\graphicspath{{figures/}}
\usepackage{booktabs}
\usepackage{array}
\usepackage{longtable}
\usepackage{multirow}
\usepackage{xcolor}
\usepackage{tikz}
\usetikzlibrary{arrows.meta,positioning,calc,shapes.geometric,shapes.misc,
                fit,backgrounds,decorations.pathmorphing}
\usepackage[most]{tcolorbox}
\usepackage[font=small,labelfont=bf,width=0.95\textwidth]{caption}
\usepackage{fancyhdr}
\usepackage{enumitem}
\usepackage{titlesec}
\usepackage{float}        % figure[H] — ตรึงรูปไว้ข้างเนื้อความที่อธิบายมัน
\usepackage{microtype}
\usepackage{listings}
\usepackage[colorlinks=true,allcolors=ctr,breaklinks=true,
            bookmarksnumbered=true]{hyperref}

% ภาษาไทยต้องการระยะบรรทัดมากกว่าละติน สระบนซ้อนวรรณยุกต์ และมีสระล่างด้วย
% ที่ระยะปกติจะชนกันข้ามบรรทัด — 1.28 คือค่าที่ทดสอบแล้ว
\linespread{1.28}
\setlength{\parskip}{0.35em}
\setlength{\parindent}{0pt}

% ---------------------------------------------------------------------
%  สี — ระบบสีเชิงความหมาย (semantic colour)
%
%  ฐานสีคือชุด Okabe–Ito ซึ่งปลอดภัยสำหรับผู้ที่มองสีบางคู่ไม่แยก (colour-blind safe)
%
%  สีสี่สีหลักถูก "ตั้งชื่อกลาง ๆ" ไว้ที่นี่ แล้วให้แต่ละเอกสารประกาศความหมายของ
%  ตัวเองในหน้า "เอกสารนี้อ่านอย่างไร":
%
%    trpc-guide.tex   → 4 ชั้นของระบบ:   สัญญา / backend / frontend / เครือข่าย
%    trpc-meeting.tex → 4 ระดับเร่งด่วน: P0 / P1 / P2 / P3
%
%  สองเอกสารใช้คนละความหมายได้ เพราะแต่ละฉบับ "ประกาศระบบสีของตัวเองไว้หน้าแรก"
%  และคงเส้นคงวาภายในเล่ม — สิ่งที่ห้ามคือเปลี่ยนความหมายกลางเล่ม
% ---------------------------------------------------------------------
\definecolor{ctr}{HTML}{0072B2}      % ฟ้า   — สัญญา/schema/type   · หรือ P2
\definecolor{beo}{HTML}{D55E00}      % ส้ม   — ฝั่ง backend         · หรือ P1
\definecolor{feg}{HTML}{009E73}      % เขียว — ฝั่ง frontend        · หรือ P3
\definecolor{ops}{HTML}{7B5EA7}      % ม่วง  — เครือข่าย/รันไทม์/ops

\definecolor{ink}{HTML}{1A2733}      % เข้ม — หัวข้อรอง
\definecolor{muted}{HTML}{5B6B7A}    % ข้อความรอง เส้นคั่น ป้ายกำกับ
\definecolor{rule}{HTML}{C9D2DB}     % เส้นบาง ขอบไดอะแกรม
\definecolor{wash}{HTML}{F1F6FA}     % พื้นกล่องฟ้า/กลาง
\definecolor{warnwash}{HTML}{FDF3E7} % พื้นกล่องส้ม
\definecolor{goodwash}{HTML}{E9F4F0} % พื้นกล่องเขียว
\definecolor{badwash}{HTML}{FBEDED}  % พื้นกล่องแดง
\definecolor{bad}{HTML}{B3261E}      % แดง — ผิด/อันตราย · หรือ P0
\definecolor{codebg}{HTML}{F7F9FB}   % พื้นบล็อกโค้ด

% ---------------------------------------------------------------------
%  HEADINGS
%  \part ใช้เลขอารบิก เพื่อให้อ้างอิงข้ามเล่มอ่านว่า "ภาค 3" ไม่ใช่ "ภาค III"
% ---------------------------------------------------------------------
\renewcommand{\thepart}{\arabic{part}}
\titleformat{\part}[display]
  {\sansall\huge\bfseries\color{ctr}}
  {\normalsize\color{muted}ภาค \thepart}{0.4em}{}
\titlespacing*{\part}{0pt}{0pt}{1.5em}

\titleformat{\section}
  {\sansall\Large\bfseries\color{ctr!88!black}}{\thesection}{0.7em}{}
\titleformat{\subsection}
  {\sansall\large\bfseries\color{ink}}{\thesubsection}{0.6em}{}
\titleformat{\subsubsection}
  {\sansall\normalsize\bfseries\color{muted}}{\thesubsubsection}{0.55em}{}
\titlespacing*{\section}{0pt}{1.9em}{0.7em}
\titlespacing*{\subsection}{0pt}{1.3em}{0.45em}
\titlespacing*{\subsubsection}{0pt}{1.0em}{0.35em}

% หัวกระดาษ: ชื่อเอกสารคงที่ทางซ้าย ชื่อหัวข้อปัจจุบันทางขวา ผู้อ่านจะรู้ตลอดว่าอยู่ตรงไหน
% \docshort ถูกกำหนดในไฟล์เอกสารแต่ละฉบับ
\providecommand{\docshort}{}
\pagestyle{fancy}
\fancyhf{}
\renewcommand{\headrulewidth}{0.4pt}
\fancyhead[L]{\footnotesize\sansall\color{muted}\docshort}
\fancyhead[R]{\footnotesize\sansall\color{muted}\leftmark}
\fancyfoot[C]{\footnotesize\color{muted}\thepage}
\renewcommand{\sectionmark}[1]{\markboth{#1}{}}

% รายการแบบกระชับ — ค่าเริ่มต้นของ LaTeX ดูหลวมเมื่ออยู่กับ \parskip
\setlist[itemize]{itemsep=1pt,topsep=3pt,parsep=0pt,leftmargin=1.4em}
\setlist[enumerate]{itemsep=1pt,topsep=3pt,parsep=0pt,leftmargin=1.6em}
\setlist[description]{itemsep=2pt,topsep=3pt,parsep=0pt}

% ---------------------------------------------------------------------
%  กล่องข้อความ 4 แบบ ความหมายตายตัว ห้ามคิดแบบที่ห้ากลางเล่ม
%
%  colbacktitle ต้องตั้งให้เท่ากับ colback มิฉะนั้น tcolorbox จะใช้ colframe
%  (สีเข้ม) เป็นพื้นหลังของหัวกล่อง แล้วตัวอักษรหัวกล่องสีเข้มจะอ่านไม่ออก
% ---------------------------------------------------------------------
\tcbset{guidebox/.style={
  enhanced, breakable, boxrule=0pt, leftrule=2.5pt, arc=1pt,
  left=8pt, right=8pt, top=6pt, bottom=6pt,
  fonttitle=\sansall\bfseries\small, titlerule=0pt,
  toptitle=1pt, bottomtitle=1pt}}

% ฟ้า — ใจความที่ต้องจำ / ประโยคที่ร้อยหลายบทเข้าด้วยกัน
\newtcolorbox{keybox}[1][]{guidebox, colback=wash, colframe=ctr,
  colbacktitle=wash, coltitle=ctr!85!black, #1}
% ส้ม — กับดัก หรือข้อสมมติที่ถ้าลืมแล้วข้อสรุปจะผิด
\newtcolorbox{warnbox}[1][]{guidebox, colback=warnwash, colframe=beo,
  colbacktitle=warnwash, coltitle=beo!80!black, #1}
% เขียว — สิ่งที่ตรวจสอบกับของจริงในรีโปแล้ว
\newtcolorbox{goodbox}[1][]{guidebox, colback=goodwash, colframe=feg,
  colbacktitle=goodwash, coltitle=feg!65!black, #1}
% แดง — วิธีที่ผิดจริง ๆ อย่าทำ
\newtcolorbox{badbox}[1][]{guidebox, colback=badwash, colframe=bad,
  colbacktitle=badwash, coltitle=bad, #1}

% ---------------------------------------------------------------------
%  SEMANTIC MARKUP
% ---------------------------------------------------------------------
% ศัพท์อังกฤษที่แทรกในประโยคภาษาไทย ฟอนต์เนื้อความครอบคลุมละตินอยู่แล้วจึงไม่ต้อง
% สลับฟอนต์ — เก็บแมโครไว้เป็น "เครื่องหมายเชิงความหมาย" เผื่อวันหลังอยากเปลี่ยน
% การแสดงผลทั้งเล่มพร้อมกัน
\newcommand{\en}[1]{#1}
% ชื่อไฟล์ ฟังก์ชัน หรือค่าที่ปรากฏในโค้ดจริง
% ไม่บังคับขนาดตัวอักษร ปล่อยให้รับขนาดจากบริบท (สำคัญในตารางที่ใช้ \footnotesize
% มิฉะนั้นโค้ดจะกลายเป็นตัวใหญ่กว่าข้อความรอบข้างแล้วดันคอลัมน์ล้น)
% ฟอนต์ mono ถูกย่อไว้ 0.84 แล้วใน \setmonofont จึงพอดีกับเนื้อความอยู่แล้ว
\newcommand{\code}[1]{{\ttfamily #1}}

% ชื่อไฟล์ที่มีอักษรไทยอยู่ในชื่อ — ห้ามใช้ \code เพราะ DejaVu Sans Mono ไม่มี
% สัญลักษณ์ไทย ตัวอักษรจะหายไปจาก PDF โดยขึ้นแค่ warning "Missing character"
% ใช้ Garuda (sans ที่มีทั้งไทยและละติน) แทน
\newcommand{\fnth}[1]{{\sffamily\small #1}}

% ป้ายชื่อแนวคิด — สีเดียวกับที่ใช้ในไดอะแกรมและตาราง
\newcommand{\stg}[2]{\textcolor{#1}{\textbf{\en{#2}}}}

% ---------------------------------------------------------------------
%  ไวยากรณ์ของไดอะแกรม — นิยามครั้งเดียว ทุกรูปจะดูเป็นชุดเดียวกัน
%  หมายเหตุ: node ที่มีข้อความไทยต้องเผื่อ inner ysep มากกว่าไดอะแกรมภาษาละติน
%  เพราะสระบน/ล่างกินที่แนวตั้ง
% ---------------------------------------------------------------------
\tikzset{
  blk/.style={draw=rule, line width=0.6pt, rounded corners=2.5pt,
              fill=wash, minimum height=9mm, inner xsep=6pt, inner ysep=5pt,
              align=center, font=\sansall\footnotesize},
  stage/.style={blk, fill=ctr!8, draw=ctr!45},
  bestage/.style={blk, fill=beo!8, draw=beo!45},
  festage/.style={blk, fill=feg!8, draw=feg!45},
  opstage/.style={blk, fill=ops!8, draw=ops!45},
  ar/.style={-{Stealth[length=2.2mm,width=1.8mm]}, line width=0.7pt,
             draw=muted},
  lbl/.style={font=\sansall\scriptsize, text=muted, align=center},
  tag/.style={font=\sansall\scriptsize\bfseries, align=center},
}

% ---------------------------------------------------------------------
%  ตาราง
%  คอลัมน์แบบชิดซ้าย: ข้อความจัดชิดขอบสองข้าง (justified) ในคอลัมน์แคบจะเปิด
%  ช่องไฟกว้างมาก และแย่เป็นพิเศษกับภาษาไทย เพราะ ICU ตัดไทยเป็นก้อนที่ใหญ่กว่า
%  คำภาษาอังกฤษ → ใช้ L{} แทน p{} ในทุกคอลัมน์ที่มีข้อความยาว
% ---------------------------------------------------------------------
\newcolumntype{L}[1]{>{\raggedright\arraybackslash}p{#1}}
\captionsetup{skip=6pt}

% ---------------------------------------------------------------------
%  บล็อกโค้ด
%
%  ข้อจำกัดที่ต้องรู้: listings อ่านไฟล์แบบไบต์ ไม่เข้าใจ UTF-8 หลายไบต์
%  → "ห้ามใส่ภาษาไทยในบล็อกโค้ด" คอมเมนต์ในโค้ดทุกอันเป็นภาษาอังกฤษ
%    คำอธิบายภาษาไทยให้อยู่นอกบล็อกเสมอ
% ---------------------------------------------------------------------
\lstdefinelanguage{TypeScript}{
  keywords={await, async, break, case, catch, class, const, continue, default,
            delete, do, else, enum, export, extends, false, finally, for, from,
            function, if, implements, import, in, instanceof, interface, let,
            new, null, of, private, public, readonly, return, super, switch,
            this, throw, true, try, type, typeof, undefined, var, void, while,
            yield, satisfies, as},
  keywordstyle=\color{ctr}\bfseries,
  ndkeywords={string, number, boolean, Date, Promise, z, Record, Array},
  ndkeywordstyle=\color{ops},
  sensitive=true,
  comment=[l]{//},
  morecomment=[s]{/*}{*/},
  commentstyle=\color{muted}\itshape,
  stringstyle=\color{feg!70!black},
  morestring=[b]',
  morestring=[b]",
  morestring=[b]`,
}

\lstset{
  language=TypeScript,
  basicstyle=\ttfamily\scriptsize,
  backgroundcolor=\color{codebg},
  frame=leftline,
  framerule=2pt,
  rulecolor=\color{rule},
  framesep=7pt,
  xleftmargin=10pt,
  breaklines=true,
  breakatwhitespace=true,
  columns=fullflexible,
  keepspaces=true,
  showstringspaces=false,
  aboveskip=0.9em,
  belowskip=0.9em,
  literate={->}{{$\rightarrow$}}2 {=>}{{=>}}2,
}

% ชื่อ float ภาษาไทย — ถ้าไม่ตั้ง จะได้ "Figure 3" อยู่ใต้เนื้อความภาษาไทย
\renewcommand{\figurename}{รูปที่}
\renewcommand{\tablename}{ตารางที่}
\renewcommand{\refname}{เอกสารอ้างอิง}
\renewcommand{\contentsname}{สารบัญ}
% หมายเหตุ: ห้ามใส่ \color ใน \contentsname เพราะ \tableofcontents ส่งมันผ่าน
% \MakeUppercase ซึ่งจะเปลี่ยน ctr เป็น CTR แล้ว xcolor จะฟ้องว่าไม่รู้จักสีนั้น
% ถ้าจะจัดสไตล์หัวสารบัญ ให้ทำผ่าน \titleformat{\section} แทน
:
%    1) \hypersetup{pdftitle={...}}
%    2) \begin{document}
%
%  Build:  cd docs && latexmk trpc-guide.tex     (.latexmkrc บังคับ xelatex ให้แล้ว)
%          PDF ออกที่ docs/build/ แล้วค่อยคัดลอกไป docs/report/
% =============================================================================

\documentclass[11pt,a4paper]{article}

% ---------------------------------------------------------------------
%  FONTS
%
%  ต้องใช้ XeLaTeX เท่านั้น — pdflatex เรียงพิมพ์ภาษาไทยไม่ได้เลย (ไม่ใช่ "ได้แต่ไม่สวย")
%
%  ใช้ตระกูล TLWG ที่มีทั้งอักษรไทยและละตินอยู่ในไฟล์เดียว จงใจเลี่ยงการจับคู่
%  ฟอนต์ละตินกับฟอนต์ไทยล้วนผ่าน ucharclasses เพราะ ucharclasses สลับฟอนต์ตาม
%  "บล็อกยูนิโคด" ของตัวอักษร โดยไม่รู้จักขอบเขตของ TeX group → ข้อความละตินที่
%  ตามหลัง \textbf{...} จะยังค้างอยู่ในฟอนต์ไทยซึ่งไม่มีสัญลักษณ์ละติน แล้ว
%  "หายไปจาก PDF เงียบ ๆ" โดยไม่มี warning
% ---------------------------------------------------------------------
\usepackage{fontspec}

\setmainfont{Laksaman}[Ligatures=TeX]          % ไทย+ละติน — เนื้อความ
\setsansfont{Garuda}[Ligatures=TeX]            % ไทย+ละติน หนักกว่า — หัวข้อ/ป้ายในไดอะแกรม
\setmonofont{DejaVu Sans Mono}[Scale=0.84]     % 0.84 ทำให้ x-height ใกล้เคียง TLWG

% --- การตัดบรรทัดภาษาไทย ---
% ภาษาไทยเขียนติดกันไม่มีช่องว่าง TeX จึงหาจุดตัดบรรทัดเองไม่ได้ และจะปล่อยให้
% บรรทัดล้นออกนอกหน้ากระดาษ สองบรรทัดนี้ยกงานให้ ICU ตัดคำแบบพจนานุกรม
% ค่า stretch ต้องมากกว่าศูนย์อย่างมีนัยสำคัญ ไม่งั้น TeX แทบไม่มีอิสระในการจัด
% บรรทัดไทยและจะออกมาเป็น overfull box แทนที่จะยืด/หดช่องไฟ
\XeTeXlinebreaklocale "th"
\XeTeXlinebreakskip = 0pt plus 0.15em minus 0.02em
\emergencystretch=2.2em

% หัวข้อ ป้ายในไดอะแกรม หัวตาราง header/footer ใช้ sans — เนื้อความไม่ใช้
% นี่คือสิ่งที่แยก "โครงสร้างสำหรับกวาดสายตา" ออกจาก "เนื้อหาสำหรับอ่าน"
\newcommand{\sansall}{\sffamily}

% ---------------------------------------------------------------------
%  LAYOUT
% ---------------------------------------------------------------------
\usepackage[a4paper,margin=2.3cm,top=2.5cm,bottom=2.5cm,headsep=12pt]{geometry}
\usepackage{amsmath,amssymb}
\usepackage{graphicx}
\graphicspath{{figures/}}
\usepackage{booktabs}
\usepackage{array}
\usepackage{longtable}
\usepackage{multirow}
\usepackage{xcolor}
\usepackage{tikz}
\usetikzlibrary{arrows.meta,positioning,calc,shapes.geometric,shapes.misc,
                fit,backgrounds,decorations.pathmorphing}
\usepackage[most]{tcolorbox}
\usepackage[font=small,labelfont=bf,width=0.95\textwidth]{caption}
\usepackage{fancyhdr}
\usepackage{enumitem}
\usepackage{titlesec}
\usepackage{float}        % figure[H] — ตรึงรูปไว้ข้างเนื้อความที่อธิบายมัน
\usepackage{microtype}
\usepackage{listings}
\usepackage[colorlinks=true,allcolors=ctr,breaklinks=true,
            bookmarksnumbered=true]{hyperref}

% ภาษาไทยต้องการระยะบรรทัดมากกว่าละติน สระบนซ้อนวรรณยุกต์ และมีสระล่างด้วย
% ที่ระยะปกติจะชนกันข้ามบรรทัด — 1.28 คือค่าที่ทดสอบแล้ว
\linespread{1.28}
\setlength{\parskip}{0.35em}
\setlength{\parindent}{0pt}

% ---------------------------------------------------------------------
%  สี — ระบบสีเชิงความหมาย (semantic colour)
%
%  ฐานสีคือชุด Okabe–Ito ซึ่งปลอดภัยสำหรับผู้ที่มองสีบางคู่ไม่แยก (colour-blind safe)
%
%  สีสี่สีหลักถูก "ตั้งชื่อกลาง ๆ" ไว้ที่นี่ แล้วให้แต่ละเอกสารประกาศความหมายของ
%  ตัวเองในหน้า "เอกสารนี้อ่านอย่างไร":
%
%    trpc-guide.tex   → 4 ชั้นของระบบ:   สัญญา / backend / frontend / เครือข่าย
%    trpc-meeting.tex → 4 ระดับเร่งด่วน: P0 / P1 / P2 / P3
%
%  สองเอกสารใช้คนละความหมายได้ เพราะแต่ละฉบับ "ประกาศระบบสีของตัวเองไว้หน้าแรก"
%  และคงเส้นคงวาภายในเล่ม — สิ่งที่ห้ามคือเปลี่ยนความหมายกลางเล่ม
% ---------------------------------------------------------------------
\definecolor{ctr}{HTML}{0072B2}      % ฟ้า   — สัญญา/schema/type   · หรือ P2
\definecolor{beo}{HTML}{D55E00}      % ส้ม   — ฝั่ง backend         · หรือ P1
\definecolor{feg}{HTML}{009E73}      % เขียว — ฝั่ง frontend        · หรือ P3
\definecolor{ops}{HTML}{7B5EA7}      % ม่วง  — เครือข่าย/รันไทม์/ops

\definecolor{ink}{HTML}{1A2733}      % เข้ม — หัวข้อรอง
\definecolor{muted}{HTML}{5B6B7A}    % ข้อความรอง เส้นคั่น ป้ายกำกับ
\definecolor{rule}{HTML}{C9D2DB}     % เส้นบาง ขอบไดอะแกรม
\definecolor{wash}{HTML}{F1F6FA}     % พื้นกล่องฟ้า/กลาง
\definecolor{warnwash}{HTML}{FDF3E7} % พื้นกล่องส้ม
\definecolor{goodwash}{HTML}{E9F4F0} % พื้นกล่องเขียว
\definecolor{badwash}{HTML}{FBEDED}  % พื้นกล่องแดง
\definecolor{bad}{HTML}{B3261E}      % แดง — ผิด/อันตราย · หรือ P0
\definecolor{codebg}{HTML}{F7F9FB}   % พื้นบล็อกโค้ด

% ---------------------------------------------------------------------
%  HEADINGS
%  \part ใช้เลขอารบิก เพื่อให้อ้างอิงข้ามเล่มอ่านว่า "ภาค 3" ไม่ใช่ "ภาค III"
% ---------------------------------------------------------------------
\renewcommand{\thepart}{\arabic{part}}
\titleformat{\part}[display]
  {\sansall\huge\bfseries\color{ctr}}
  {\normalsize\color{muted}ภาค \thepart}{0.4em}{}
\titlespacing*{\part}{0pt}{0pt}{1.5em}

\titleformat{\section}
  {\sansall\Large\bfseries\color{ctr!88!black}}{\thesection}{0.7em}{}
\titleformat{\subsection}
  {\sansall\large\bfseries\color{ink}}{\thesubsection}{0.6em}{}
\titleformat{\subsubsection}
  {\sansall\normalsize\bfseries\color{muted}}{\thesubsubsection}{0.55em}{}
\titlespacing*{\section}{0pt}{1.9em}{0.7em}
\titlespacing*{\subsection}{0pt}{1.3em}{0.45em}
\titlespacing*{\subsubsection}{0pt}{1.0em}{0.35em}

% หัวกระดาษ: ชื่อเอกสารคงที่ทางซ้าย ชื่อหัวข้อปัจจุบันทางขวา ผู้อ่านจะรู้ตลอดว่าอยู่ตรงไหน
% \docshort ถูกกำหนดในไฟล์เอกสารแต่ละฉบับ
\providecommand{\docshort}{}
\pagestyle{fancy}
\fancyhf{}
\renewcommand{\headrulewidth}{0.4pt}
\fancyhead[L]{\footnotesize\sansall\color{muted}\docshort}
\fancyhead[R]{\footnotesize\sansall\color{muted}\leftmark}
\fancyfoot[C]{\footnotesize\color{muted}\thepage}
\renewcommand{\sectionmark}[1]{\markboth{#1}{}}

% รายการแบบกระชับ — ค่าเริ่มต้นของ LaTeX ดูหลวมเมื่ออยู่กับ \parskip
\setlist[itemize]{itemsep=1pt,topsep=3pt,parsep=0pt,leftmargin=1.4em}
\setlist[enumerate]{itemsep=1pt,topsep=3pt,parsep=0pt,leftmargin=1.6em}
\setlist[description]{itemsep=2pt,topsep=3pt,parsep=0pt}

% ---------------------------------------------------------------------
%  กล่องข้อความ 4 แบบ ความหมายตายตัว ห้ามคิดแบบที่ห้ากลางเล่ม
%
%  colbacktitle ต้องตั้งให้เท่ากับ colback มิฉะนั้น tcolorbox จะใช้ colframe
%  (สีเข้ม) เป็นพื้นหลังของหัวกล่อง แล้วตัวอักษรหัวกล่องสีเข้มจะอ่านไม่ออก
% ---------------------------------------------------------------------
\tcbset{guidebox/.style={
  enhanced, breakable, boxrule=0pt, leftrule=2.5pt, arc=1pt,
  left=8pt, right=8pt, top=6pt, bottom=6pt,
  fonttitle=\sansall\bfseries\small, titlerule=0pt,
  toptitle=1pt, bottomtitle=1pt}}

% ฟ้า — ใจความที่ต้องจำ / ประโยคที่ร้อยหลายบทเข้าด้วยกัน
\newtcolorbox{keybox}[1][]{guidebox, colback=wash, colframe=ctr,
  colbacktitle=wash, coltitle=ctr!85!black, #1}
% ส้ม — กับดัก หรือข้อสมมติที่ถ้าลืมแล้วข้อสรุปจะผิด
\newtcolorbox{warnbox}[1][]{guidebox, colback=warnwash, colframe=beo,
  colbacktitle=warnwash, coltitle=beo!80!black, #1}
% เขียว — สิ่งที่ตรวจสอบกับของจริงในรีโปแล้ว
\newtcolorbox{goodbox}[1][]{guidebox, colback=goodwash, colframe=feg,
  colbacktitle=goodwash, coltitle=feg!65!black, #1}
% แดง — วิธีที่ผิดจริง ๆ อย่าทำ
\newtcolorbox{badbox}[1][]{guidebox, colback=badwash, colframe=bad,
  colbacktitle=badwash, coltitle=bad, #1}

% ---------------------------------------------------------------------
%  SEMANTIC MARKUP
% ---------------------------------------------------------------------
% ศัพท์อังกฤษที่แทรกในประโยคภาษาไทย ฟอนต์เนื้อความครอบคลุมละตินอยู่แล้วจึงไม่ต้อง
% สลับฟอนต์ — เก็บแมโครไว้เป็น "เครื่องหมายเชิงความหมาย" เผื่อวันหลังอยากเปลี่ยน
% การแสดงผลทั้งเล่มพร้อมกัน
\newcommand{\en}[1]{#1}
% ชื่อไฟล์ ฟังก์ชัน หรือค่าที่ปรากฏในโค้ดจริง
% ไม่บังคับขนาดตัวอักษร ปล่อยให้รับขนาดจากบริบท (สำคัญในตารางที่ใช้ \footnotesize
% มิฉะนั้นโค้ดจะกลายเป็นตัวใหญ่กว่าข้อความรอบข้างแล้วดันคอลัมน์ล้น)
% ฟอนต์ mono ถูกย่อไว้ 0.84 แล้วใน \setmonofont จึงพอดีกับเนื้อความอยู่แล้ว
\newcommand{\code}[1]{{\ttfamily #1}}

% ชื่อไฟล์ที่มีอักษรไทยอยู่ในชื่อ — ห้ามใช้ \code เพราะ DejaVu Sans Mono ไม่มี
% สัญลักษณ์ไทย ตัวอักษรจะหายไปจาก PDF โดยขึ้นแค่ warning "Missing character"
% ใช้ Garuda (sans ที่มีทั้งไทยและละติน) แทน
\newcommand{\fnth}[1]{{\sffamily\small #1}}

% ป้ายชื่อแนวคิด — สีเดียวกับที่ใช้ในไดอะแกรมและตาราง
\newcommand{\stg}[2]{\textcolor{#1}{\textbf{\en{#2}}}}

% ---------------------------------------------------------------------
%  ไวยากรณ์ของไดอะแกรม — นิยามครั้งเดียว ทุกรูปจะดูเป็นชุดเดียวกัน
%  หมายเหตุ: node ที่มีข้อความไทยต้องเผื่อ inner ysep มากกว่าไดอะแกรมภาษาละติน
%  เพราะสระบน/ล่างกินที่แนวตั้ง
% ---------------------------------------------------------------------
\tikzset{
  blk/.style={draw=rule, line width=0.6pt, rounded corners=2.5pt,
              fill=wash, minimum height=9mm, inner xsep=6pt, inner ysep=5pt,
              align=center, font=\sansall\footnotesize},
  stage/.style={blk, fill=ctr!8, draw=ctr!45},
  bestage/.style={blk, fill=beo!8, draw=beo!45},
  festage/.style={blk, fill=feg!8, draw=feg!45},
  opstage/.style={blk, fill=ops!8, draw=ops!45},
  ar/.style={-{Stealth[length=2.2mm,width=1.8mm]}, line width=0.7pt,
             draw=muted},
  lbl/.style={font=\sansall\scriptsize, text=muted, align=center},
  tag/.style={font=\sansall\scriptsize\bfseries, align=center},
}

% ---------------------------------------------------------------------
%  ตาราง
%  คอลัมน์แบบชิดซ้าย: ข้อความจัดชิดขอบสองข้าง (justified) ในคอลัมน์แคบจะเปิด
%  ช่องไฟกว้างมาก และแย่เป็นพิเศษกับภาษาไทย เพราะ ICU ตัดไทยเป็นก้อนที่ใหญ่กว่า
%  คำภาษาอังกฤษ → ใช้ L{} แทน p{} ในทุกคอลัมน์ที่มีข้อความยาว
% ---------------------------------------------------------------------
\newcolumntype{L}[1]{>{\raggedright\arraybackslash}p{#1}}
\captionsetup{skip=6pt}

% ---------------------------------------------------------------------
%  บล็อกโค้ด
%
%  ข้อจำกัดที่ต้องรู้: listings อ่านไฟล์แบบไบต์ ไม่เข้าใจ UTF-8 หลายไบต์
%  → "ห้ามใส่ภาษาไทยในบล็อกโค้ด" คอมเมนต์ในโค้ดทุกอันเป็นภาษาอังกฤษ
%    คำอธิบายภาษาไทยให้อยู่นอกบล็อกเสมอ
% ---------------------------------------------------------------------
\lstdefinelanguage{TypeScript}{
  keywords={await, async, break, case, catch, class, const, continue, default,
            delete, do, else, enum, export, extends, false, finally, for, from,
            function, if, implements, import, in, instanceof, interface, let,
            new, null, of, private, public, readonly, return, super, switch,
            this, throw, true, try, type, typeof, undefined, var, void, while,
            yield, satisfies, as},
  keywordstyle=\color{ctr}\bfseries,
  ndkeywords={string, number, boolean, Date, Promise, z, Record, Array},
  ndkeywordstyle=\color{ops},
  sensitive=true,
  comment=[l]{//},
  morecomment=[s]{/*}{*/},
  commentstyle=\color{muted}\itshape,
  stringstyle=\color{feg!70!black},
  morestring=[b]',
  morestring=[b]",
  morestring=[b]`,
}

\lstset{
  language=TypeScript,
  basicstyle=\ttfamily\scriptsize,
  backgroundcolor=\color{codebg},
  frame=leftline,
  framerule=2pt,
  rulecolor=\color{rule},
  framesep=7pt,
  xleftmargin=10pt,
  breaklines=true,
  breakatwhitespace=true,
  columns=fullflexible,
  keepspaces=true,
  showstringspaces=false,
  aboveskip=0.9em,
  belowskip=0.9em,
  literate={->}{{$\rightarrow$}}2 {=>}{{=>}}2,
}

% ชื่อ float ภาษาไทย — ถ้าไม่ตั้ง จะได้ "Figure 3" อยู่ใต้เนื้อความภาษาไทย
\renewcommand{\figurename}{รูปที่}
\renewcommand{\tablename}{ตารางที่}
\renewcommand{\refname}{เอกสารอ้างอิง}
\renewcommand{\contentsname}{สารบัญ}
% หมายเหตุ: ห้ามใส่ \color ใน \contentsname เพราะ \tableofcontents ส่งมันผ่าน
% \MakeUppercase ซึ่งจะเปลี่ยน ctr เป็น CTR แล้ว xcolor จะฟ้องว่าไม่รู้จักสีนั้น
% ถ้าจะจัดสไตล์หัวสารบัญ ให้ทำผ่าน \titleformat{\section} แทน


\renewcommand{\docshort}{สัญญา tRPC · วาระประชุม}
\hypersetup{pdftitle={วาระประชุมตกลงสัญญา tRPC — ULMs},
            pdfauthor={ทีมพัฒนา ULMs}}

% -----------------------------------------------------------------------------
%  ป้ายระดับความเร่งด่วน — ระบบสีเฉพาะของเอกสารฉบับนี้
%  (ฉบับ trpc-guide.tex ใช้สีชุดเดียวกันแทน "ชั้นของระบบ" ซึ่งเป็นคนละความหมาย
%   แต่ละฉบับประกาศระบบสีของตัวเองไว้ในหน้า "อ่านอย่างไร" และคงเส้นคงวาภายในเล่ม)
% -----------------------------------------------------------------------------
\newcommand{\badge}[2]{{\sansall\bfseries\scriptsize%
  \colorbox{#1!16}{\textcolor{#1!72!black}{\,#2\,}}}}
\newcommand{\pzero}{\badge{bad}{P0}}
\newcommand{\pone}{\badge{beo}{P1}}
\newcommand{\ptwo}{\badge{ctr}{P2}}
\newcommand{\pthree}{\badge{feg}{P3}}

% กล่อง "ผลลัพธ์ที่ต้องได้" — ปิดท้ายทุกวาระ เพื่อไม่ให้คุยจบแบบไม่มีมติ
\newtcolorbox{outbox}{guidebox, colback=wash, colframe=ctr,
  colbacktitle=wash, coltitle=ctr!85!black,
  title={ผลลัพธ์ที่ต้องจดลงบันทึกก่อนปิดวาระนี้}}

\begin{document}

% =============================================================================
%  หน้าปก
% =============================================================================
\begin{titlepage}
\thispagestyle{empty}
\begin{tikzpicture}[remember picture, overlay]
  \fill[ctr!8] (current page.north west) rectangle
       ([yshift=-9.5cm]current page.north east);
  \fill[ctr] ([yshift=-9.5cm]current page.north west) rectangle
       ([yshift=-9.62cm]current page.north east);
\end{tikzpicture}

\vspace*{1.4cm}
{\sansall\color{muted}\small เอกสารสำหรับใช้ในห้องประชุม --- ผู้เข้าร่วมทุกคนควรอ่านภาค 1 มาก่อน}

\vspace{0.9cm}
{\sansall\bfseries\fontsize{27}{33}\selectfont\color{ctr!90!black}
วาระประชุม\\ตกลงสัญญา \en{tRPC}\par}

\vspace{0.5cm}
{\sansall\fontsize{13.5}{19}\selectfont\color{ink}
19 เรื่องที่ต้องตกลง เรียงตามว่าอะไรบล็อกอะไร\\
เรื่องไหนต้องจบก่อนออกจากห้อง เรื่องไหนผลัดไปก่อนได้\par}

\vspace{2.8cm}

\begin{minipage}{0.66\textwidth}
เอกสารนี้ไม่ใช่บันทึกการประชุม แต่เป็น \textbf{สิ่งที่ต้องเตรียมก่อนประชุม}
แต่ละวาระมาพร้อมคำถามที่ต้องตอบ ตัวเลือกที่มี ข้อเสนอของฝ่ายเทคนิค
และผลที่จะตามมาถ้าไม่ตกลงในรอบนี้

\vspace{0.6em}
จุดประสงค์คือให้ทุกคนเข้าห้องประชุมโดยรู้ตรงกันว่า \emph{กำลังจะตัดสินใจอะไร}
และ \emph{ทำไมลำดับถึงเป็นแบบนี้} --- ไม่ใช่มาถกกันว่าเรื่องไหนสำคัญกว่ากัน

\vspace{0.6em}
เหตุผลเบื้องหลังของแต่ละวาระอยู่ในเอกสารคู่กัน \code{trpc-guide.tex}
ส่วนภาคผนวก ก ของเล่มนี้คือแบบฟอร์มบันทึกมติที่กรอกได้ระหว่างประชุม
\end{minipage}

\vfill
\begin{tikzpicture}
  \draw[rule, line width=0.6pt] (0,0) -- (\textwidth,0);
\end{tikzpicture}
\vspace{0.3em}

{\sansall\footnotesize\color{muted}
โครงการ: ระบบบริหารจัดการการยืม--คืนอุปกรณ์มหาวิทยาลัย (ULMs) \\
เวลาที่เสนอ: 90 นาที สำหรับภาค 1--2 (\pzero{} และ \pone{}) · ภาค 3--4 ไม่ต้องคุยในรอบนี้ \\
ปรับปรุงล่าสุด: 8 สิงหาคม 2569}
\end{titlepage}

% =============================================================================
%  สารบัญ
% =============================================================================
\microtypesetup{protrusion=false}
{\hypersetup{allcolors=ink}
\tableofcontents}
\microtypesetup{protrusion=true}
\clearpage

% =============================================================================
%  เอกสารนี้อ่านอย่างไร
% =============================================================================
\section*{เอกสารนี้อ่านอย่างไร}
\addcontentsline{toc}{section}{เอกสารนี้อ่านอย่างไร}
\markboth{เอกสารนี้อ่านอย่างไร}{}

เอกสารแบ่งเป็นห้าภาค \textbf{เรียงตามว่าเรื่องไหนบล็อกเรื่องไหน} ไม่ได้เรียงตาม
ความยากหรือตามลำดับที่จะเขียนโค้ด เหตุผลคือในการประชุมที่มีเวลาจำกัด สิ่งที่ทำให้
ประชุมล้มเหลวไม่ใช่การตกลงไม่ได้ แต่คือ \textbf{การใช้เวลาไปกับเรื่องที่รอได้
จนเรื่องที่รอไม่ได้ไม่ได้คุย}

\subsection*{ระบบสีในเอกสารฉบับนี้ --- สี่ระดับความเร่งด่วน}

\begin{table}[H]
\centering\small
\begin{tabular}{L{1.3cm}L{4.6cm}L{9.0cm}}
\toprule
\sansall\bfseries ระดับ & \sansall\bfseries กำหนด & \sansall\bfseries เกณฑ์ที่ใช้จัดระดับ \\
\midrule
\pzero{} & ต้องจบก่อนออกจากห้อง &
ถ้าไม่ตกลง \textbf{ไม่มีใครเขียนโค้ดต่อได้เลย} ทั้งสองฝั่ง \\
\addlinespace
\pone{} & ต้องจบภายในสัปดาห์นี้ &
เขียน \en{procedure} ตัวแรกได้ แต่ถ้ายังไม่ตกลง \en{procedure} ตัวที่ 2--30
จะออกมาไม่เหมือนกัน แล้วต้องกลับมาแก้ทั้งหมด \\
\addlinespace
\ptwo{} & ตกลงตอนเริ่มทำ \en{feature} นั้น &
กระทบเฉพาะบางส่วนของระบบ ตัดสินใจช้าแล้วแก้ทีหลังได้ในต้นทุนที่ยอมรับได้ \\
\addlinespace
\pthree{} & ผลัดไปหลัง \en{MVP} ได้ &
ไม่กระทบการส่งงาน ถ้าคุยตอนนี้จะกินเวลาโดยไม่ได้อะไรกลับมา \\
\bottomrule
\end{tabular}
\caption{สี่ระดับ --- เกณฑ์คือ ``ถ้าไม่ตกลงตอนนี้ จะต้องกลับมาแก้ของที่ทำไปแล้วมากแค่ไหน'' ไม่ใช่ ``เรื่องนี้สำคัญแค่ไหน''}
\end{table}

\begin{warnbox}[title={ทุกวาระสำคัญ --- แต่ไม่ใช่ทุกวาระเร่งด่วน}]
วาระใน \pthree{} ไม่ได้แปลว่าไม่สำคัญ การจำกัดอัตราการเรียก (\en{rate limit})
สำคัญมากในระบบจริง แต่มันแก้ทีหลังได้โดยไม่ต้องรื้อของเดิม ต่างจากการตั้งชื่อ
\en{procedure} ที่ถ้าเปลี่ยนตอนมี 30 ตัวแล้วจะกระทบทุกไฟล์ทั้งสองฝั่ง

การจัดลำดับนี้ใช้เกณฑ์ \textbf{ต้นทุนของการตัดสินใจช้า} ไม่ใช่ความสำคัญ
\end{warnbox}

\subsection*{แผนที่ความสัมพันธ์ --- อะไรบล็อกอะไร}

รูปข้างล่างคือเหตุผลของลำดับทั้งหมด ลูกศรอ่านว่า ``ต้องตกลงข้อนี้ก่อน จึงจะตกลง
ข้อถัดไปได้อย่างมีความหมาย''

\begin{figure}[H]
\centering
\begin{tikzpicture}[node distance=4mm, every node/.style={font=\sansall\scriptsize}]

  % ---- P0 ----
  \node[blk, fill=badwash, draw=bad!45, text width=3.2cm] (d1)
       {\textbf{ว-01} ขนส่งสัญญา\\[1pt]\scriptsize\color{muted}monorepo / \mbox{คัดลอก} / paths};
  \node[blk, fill=badwash, draw=bad!45, text width=3.2cm, right=of d1] (d2)
       {\textbf{ว-02} เวอร์ชัน \code{zod}};
  \node[blk, fill=badwash, draw=bad!45, text width=3.2cm, right=of d2] (d3)
       {\textbf{ว-03} \code{ctx.user} + \en{cookie}};
  \node[blk, fill=badwash, draw=bad!45, text width=3.2cm, right=of d3] (d4)
       {\textbf{ว-04} เจ้าของสัญญา + วิธีแก้};

  % ---- P1 ----
  \node[blk, fill=warnwash, draw=beo!45, text width=3.2cm, below=11mm of d1] (d5)
       {\textbf{ว-05} ชื่อ \en{router}/\en{procedure}};
  \node[blk, fill=warnwash, draw=beo!45, text width=3.2cm, right=of d5] (d6)
       {\textbf{ว-06} รูปแบบ \en{error}};
  \node[blk, fill=warnwash, draw=beo!45, text width=3.2cm, right=of d6] (d7)
       {\textbf{ว-07} แบ่งหน้า/กรอง};
  \node[blk, fill=warnwash, draw=beo!45, text width=3.2cm, right=of d7] (d8)
       {\textbf{ว-08} รูปแบบวันที่};

  \node[blk, fill=warnwash, draw=beo!45, text width=5.2cm, below=6mm of d5, xshift=1.1cm] (d9)
       {\textbf{ว-09} กฎการแปลงชื่อฟิลด์ + ห้ามส่ง \en{Prisma model} ตรง ๆ};
  \node[blk, fill=warnwash, draw=beo!45, text width=5.2cm, right=of d9] (d10)
       {\textbf{ว-10} สถานะเป็นสตริงคงที่ ไม่ใช่เลข \en{key}};

  % ---- P2/P3 รวบ ----
  \node[blk, fill=wash, draw=ctr!45, text width=7.0cm, below=11mm of d9, xshift=1.0cm] (p2)
       {\textbf{ว-11 ถึง ว-15} อัปโหลด · ความถี่ \en{polling} · \en{batching} · ข้อความ · \en{log}\\[1pt]
        \scriptsize\color{muted}ตกลงตอนเริ่มทำ \en{feature} นั้น};
  \node[blk, fill=goodwash, draw=feg!45, text width=6.2cm, right=of p2] (p3)
       {\textbf{ว-16 ถึง ว-19} \en{realtime} · \en{rate limit} · \en{versioning} · เทสต์สัญญา\\[1pt]
        \scriptsize\color{muted}หลัง \en{MVP}};

  \draw[ar] (d1.south) -- ++(0,-3mm) -| (d5.north);
  \draw[ar] (d1.south) -- ++(0,-3mm) -| (d9.north);
  \draw[ar] (d2.south) -- ++(0,-3mm) -| ([xshift=-1cm]d5.north);
  \draw[ar] (d3.south) -- ++(0,-3mm) -| (d6.north);
  \draw[ar] (d4.south) -- ++(0,-3mm) -| (d8.north);
  \draw[ar] (d9.south) -- ++(0,-3mm) -| ([xshift=-1.5cm]p2.north);
  \draw[ar] (d10.south) -- ++(0,-3mm) -| ([xshift=1.5cm]p2.north);
  \draw[ar, dashed] (p2) -- (p3);

\end{tikzpicture}
\caption{แผนที่วาระ --- แถวบนสุดคือสี่เรื่องที่ทุกอย่างข้างล่างพึ่งพา ถ้าประชุมได้แค่ 30 นาที ให้คุยเฉพาะแถวบน ลูกศรประหมายถึง ``คุยได้แต่ไม่จำเป็น'' คือภาค 3 กับ 4 ที่ไม่ควรอยู่ในวาระของรอบนี้}
\label{fig:dep}
\end{figure}

\subsection*{แต่ละวาระเขียนในรูปแบบเดียวกัน}

เพื่อให้ประชุมเร็วและไม่หลงประเด็น ทุกวาระมีหัวข้อย่อยชุดเดียวกัน:
\textbf{คำถามที่ต้องตอบ} $\to$ \textbf{ทำไมอยู่ระดับนี้} $\to$
\textbf{ตัวเลือก} $\to$ \textbf{ข้อเสนอ} $\to$ \textbf{ถ้าไม่ตกลงจะเกิดอะไร}
$\to$ \textbf{ผลลัพธ์ที่ต้องจด}

กล่องฟ้าปิดท้ายทุกวาระคือรายการที่ต้อง \emph{จดลงบันทึกจริง} --- ถ้าคุยจบแล้ว
กรอกกล่องนั้นไม่ได้ แปลว่ายังไม่ได้ตกลง

\clearpage

% =============================================================================
\part{\pzero{} ต้องจบก่อนออกจากห้อง}
% =============================================================================

\textbf{สี่วาระนี้บล็อกทุกคน ถ้าไม่จบ ทั้ง \en{frontend} และ \en{backend}
เขียนโค้ดที่จะได้ใช้จริงไม่ได้เลย เวลาที่เสนอ: 55 นาที}

\section{ว-01 · ไฟล์สัญญาจะเดินทางข้ามฝั่งอย่างไร}

\textbf{คำถามที่ต้องตอบ:} \en{frontend} จะ \code{import type} ไฟล์สัญญาที่
\en{backend} สร้างขึ้นได้ด้วยวิธีไหน

\textbf{ทำไมอยู่ระดับ \pzero{}:} สัญญาของ \en{tRPC} คือ \en{type} ของ
\en{TypeScript} ถ้าฝั่งหนึ่งมองไม่เห็นไฟล์ของอีกฝั่ง ก็ไม่มีสัญญา ---
มีแต่ \en{REST} ที่เขียนด้วยไวยากรณ์ของ \en{tRPC} ตอนนี้รีโปมี
\code{package.json} สามไฟล์ที่ไม่รู้จักกัน --- ไม่มี \en{workspace} เชื่อม
จึง \code{import type} ข้ามกันไม่ได้ (หมายเหตุ: \code{frontend-preview/} กับ
\code{backend-preview/} เป็น \en{git worktree} ของอีกสองสาขา โค้ดจึงอยู่บนดิสก์
พร้อมกันแล้ว สคริปต์คัดลอกในทางเลือก ข. รันข้ามได้ทันทีในเครื่องนักพัฒนา)

\subsection{คำถามแรกที่จะถูกถามในห้อง: ``\en{merge} รวมสาขาให้จบก็แก้ได้ไม่ใช่เหรอ''}

คำถามนี้สมเหตุสมผลมากและต้องตอบให้จบก่อน ไม่งั้นจะวนกลับมาถามซ้ำระหว่างคุย
ทางเลือก คำตอบคือ \textbf{\en{merge} แก้ได้ครึ่งเดียว และครึ่งที่เหลือคือครึ่งที่
บล็อกเราอยู่จริง ๆ}

\begin{goodbox}[title={ข่าวดีก่อน: งาน \en{merge} เหลือน้อยกว่าที่คิด และไม่มี \en{conflict}}]
ตรวจกับ \code{origin} จริงเมื่อ 8 ส.ค. 2569 (หลังทีมทยอย \en{merge} เข้า
\code{main} ไปแล้วหลายรอบ)

\begin{itemize}
\item \code{front-end-branch} --- \textbf{ถูก \en{merge} เข้า \code{main}
      ไปเรียบร้อยแล้ว} เหลือ 0 \en{commit} ที่ค้าง
\item \code{back-end-branch} --- เหลือ 2 \en{commit} (\code{Set up Nest.js+Prisma}
      และ \code{Added tRPC}) ส่วน \code{main} นำอยู่ 19 \en{commit}
\item \textbf{\en{conflict} ที่จะเกิดจากการรวม \code{back-end-branch} เข้า
      \code{main}: 0 จุด} (ตรวจด้วย \code{git merge-tree --write-tree})
\end{itemize}

แปลว่างานที่เหลือคือ \en{merge} สาขาเดียว 2 \en{commit} แบบไม่มีอะไรชน ---
การ \en{merge} ไม่ใช่งานที่น่ากลัว และควรทำ แต่ด้วยเหตุผลอื่นที่ไม่ใช่เรื่องสัญญา
\end{goodbox}

สมมติว่า \en{merge} เสร็จแล้ว ทุกอย่างอยู่บน \code{main} โฟลเดอร์เดียวกัน
แล้วฝั่ง \en{frontend} เขียนบรรทัดนี้

\begin{lstlisting}
// after the merge — everything is in one branch, one folder tree
import type { AppRouter } from '../../backend/src/@generated/appRouter';
\end{lstlisting}

จะเจอด่านสามชั้นเรียงกัน และ \textbf{ไม่มีด่านไหนเลยที่เกี่ยวกับสาขา \en{git}}

\begin{figure}[H]
\centering
\begin{tikzpicture}[node distance=4mm]
  \node[festage, text width=3.0cm] (fe)
       {\en{frontend} เขียน\\\code{import type}};
  \node[blk, fill=warnwash, draw=beo!45, text width=3.1cm, right=of fe] (g1)
       {\textbf{ด่าน 1}\\\code{tsconfig include}\\[1pt]
        \scriptsize\color{muted}\en{TS} ไม่มองไฟล์นอก \code{src/}};
  \node[blk, fill=badwash, draw=bad!55, text width=3.4cm, right=of g1] (g2)
       {\textbf{ด่าน 2}\\\code{node\_modules} สองกอง\\[1pt]
        \scriptsize\color{muted}ด่านที่แก้ยากจริง};
  \node[blk, fill=warnwash, draw=beo!45, text width=3.0cm, right=of g2] (g3)
       {\textbf{ด่าน 3}\\\code{zod} 3 กับ 4\\[1pt]
        \scriptsize\color{muted}คือวาระ ว-02};

  \draw[ar] (fe) -- (g1); \draw[ar] (g1) -- (g2); \draw[ar] (g2) -- (g3);

  \node[blk, fill=goodwash, draw=feg!45, text width=3.1cm, below=9mm of g1] (m1)
       {\en{merge} ช่วยไหม?\\[1pt]\scriptsize\color{muted}\textbf{ไม่} แต่แก้ง่าย เพิ่มบรรทัดเดียว};
  \node[blk, fill=badwash, draw=bad!45, text width=3.4cm, below=9mm of g2] (m2)
       {\en{merge} ช่วยไหม?\\[1pt]\scriptsize\color{muted}\textbf{ไม่ช่วยเลย} เพราะตัวที่แยกคือ \code{package.json} ไม่ใช่สาขา};
  \node[blk, fill=badwash, draw=bad!45, text width=3.0cm, below=9mm of g3] (m3)
       {\en{merge} ช่วยไหม?\\[1pt]\scriptsize\color{muted}\textbf{ไม่} ต้องตกลง ว-02 อยู่ดี};

  \draw[ar, draw=rule] (g1) -- (m1);
  \draw[ar, draw=rule] (g2) -- (m2);
  \draw[ar, draw=rule] (g3) -- (m3);
\end{tikzpicture}
\caption{สามด่านที่ \code{import type} ต้องผ่าน --- แถวล่างคือคำตอบว่า \en{merge} ช่วยด่านไหนบ้าง}
\label{fig:merge-gates}
\end{figure}

\subsubsection{ด่าน 1 --- \en{TypeScript} ไม่มองไฟล์นั้นตั้งแต่แรก}

\code{frontend/tsconfig.json} ระบุ \code{"include": ["src", "vite.config.ts",
"tailwind.config.ts"]} ไฟล์ที่อยู่นอก \code{src/} จึงไม่ถูกคอมไพล์เลย

ด่านนี้ \textbf{แก้ง่ายที่สุด} แค่เติมพาธเข้าไปใน \code{include} หรือเพิ่ม
\code{paths} --- แต่พอผ่านด่านนี้ไปได้ จะไปเจอด่าน 2 ทันที และด่าน 2 คือด่านจริง

\subsubsection{ด่าน 2 --- \code{node\_modules} สองกองแยกกัน (ปัญหาของการไม่ใช้ \en{npm} เดียวกัน)}

นี่คือหัวใจของทั้งวาระ ต้องเข้าใจกลไกให้ตรงกันก่อนเลือกทางเลือก

ไฟล์สัญญาที่ \en{backend} สร้างขึ้น \textbf{ไม่ใช่ \en{type} ล้วน ๆ} มันมีโค้ด
\code{zod} อยู่ข้างในและ \code{import} ของมันเองต่อ

\begin{lstlisting}
// inside the generated contract file, which physically lives in backend/
import { z } from 'zod';
import { initTRPC } from '@trpc/server';

const userOutput = z.object({ id: z.number(), firstName: z.string() });
// ...
\end{lstlisting}

เมื่อ \en{frontend} ดึงไฟล์นี้เข้ามาในโปรแกรมของตัวเอง \en{TypeScript} ต้อง
ตามไป \en{resolve} \code{'zod'} และ \code{'@trpc/server'} ให้ได้ด้วย ---
และตรงนี้แหละที่พัง เพราะ

\begin{enumerate}
\item \textbf{\en{npm} ติดตั้งแยกกอง} --- \code{frontend/node\_modules} กับ
      \code{backend/node\_modules} เป็นคนละต้นไม้ ไม่มีการใช้ร่วมกัน
      \en{merge} สาขาไม่ได้ทำให้สองกองนี้รวมกัน เพราะสิ่งที่กำหนดว่ามีกี่กองคือ
      \textbf{จำนวนไฟล์ \code{package.json}} ไม่ใช่จำนวนสาขา
\item \textbf{\en{frontend} ไม่มี \code{@trpc/server} เลย} --- ตรวจแล้วใน
      \code{frontend/package.json} ไม่มีแพ็กเกจของ \en{tRPC} สักตัว
      \en{resolve} ไม่เจอ ผลคือ \en{type} ทั้งก้อนกลายเป็น \code{any} หรือ
      \en{error}
\item \textbf{\code{zod} คนละเวอร์ชันใหญ่} --- 3.23 กับ 4.4 (ด่าน 3)
\end{enumerate}

\begin{warnbox}[title={อาการเฉพาะที่ควรจำไว้: ``\en{Type} X ใช้แทน \en{Type} X ไม่ได้''}]
สมมติแก้ข้อ 2 และ 3 แล้ว คือ \en{frontend} ลง \code{zod} เวอร์ชันเดียวกันเป๊ะ
แต่ยัง \en{install} แยกกองอยู่ --- ก็ยังมีโอกาสพังได้อีกแบบ

\en{TypeScript} ตัดสินว่า \en{type} สองตัวเป็นตัวเดียวกันไหมจาก
\textbf{ไฟล์ประกาศที่มันมาจาก} ไม่ใช่จากชื่อ ถ้ามี \code{zod} สองชุดบนดิสก์
มันจะเห็น \code{ZodString} \textbf{สองชนิดที่ไม่เท่ากัน} แล้วฟ้อง \en{error}
หน้าตาแบบ ``\code{Type 'ZodString' is not assignable to type 'ZodString'}''
ซึ่งอ่านแล้วไม่มีทางเดาสาเหตุถูก และเป็นบั๊กที่กินเวลาทีมได้ทั้งวัน

\en{workspaces} แก้ปัญหานี้โดยตรง เพราะมันยุบให้เหลือชุดเดียวจริง ๆ
\end{warnbox}

\subsubsection{ด่าน 3 --- \code{zod} คนละเวอร์ชัน}

รายละเอียดอยู่ในวาระ ว-02 ที่ต้องเน้นตรงนี้คือ \en{merge} ไม่ได้ทำให้เวอร์ชัน
ตรงกันเอง ยังต้องมีคนตัดสินใจและลงมือ

\begin{keybox}[title={ประโยคที่ใช้ตอบในห้องประชุมได้เลย}]
``\en{merge} ทำเถอะ ง่ายมาก ไม่มี \en{conflict} เลย แต่มันไม่ทำให้
\code{import type} ข้ามฝั่งได้ เพราะตัวที่แยกคือ \code{package.json} สองไฟล์
ไม่ใช่สาขา --- เรายังต้องเลือกระหว่าง \en{workspaces} กับคัดลอกไฟล์อยู่ดี''
\end{keybox}

\subsubsection{แล้ว \en{merge} ให้อะไรเรา}

ไม่ใช่ว่าไม่คุ้ม --- \en{merge} \textbf{ควรทำ} แต่ต้องรู้ว่าได้อะไรจริง

\begin{itemize}
\item ทุกคนเห็นโค้ดชุดเดียวกัน รีวิว \en{PR} ข้ามฝั่งได้
\item \textbf{เป็นเงื่อนไขที่ต้องมีก่อน ถ้าจะไปทางเลือก ก.} เพราะ
      \code{package.json} ระดับรากต้องอยู่บนสาขาเดียวกับทั้งสองโฟลเดอร์
\item ทำให้ \en{CI} ในอนาคตตรวจทั้งสองฝั่งได้ในการรันเดียว
\item ลดโอกาสที่โค้ดสองฝั่งจะห่างกันจนรวมยากในภายหลัง
\end{itemize}

\begin{badbox}[title={ทางเลือก ค. (\code{paths} ใน \code{tsconfig}) ตายตรงด่าน 2}]
ทางเลือกที่เคยเสนอไว้ว่า ``ชี้ \code{paths} ข้ามโฟลเดอร์เอา ไม่ต้องคัดลอก''
คือการแก้เฉพาะด่าน 1 แล้วเดินชนด่าน 2 เต็ม ๆ

มันไม่ได้ทำอะไรกับ \code{node\_modules} สองกองเลย จึงไม่ใช่ทางเลือกจริง
--- \textbf{ตัดออกจากการพิจารณาได้ เหลือแค่ ก. กับ ข.}
\end{badbox}

\subsection{ทางเลือก ก. --- \en{npm workspaces}}

\subsubsection{แนวคิด}

บอก \en{npm} ว่า \code{frontend/} กับ \code{backend/} เป็น ``\en{workspace}''
ของรีโปเดียวกัน \en{npm} จะ \textbf{ยุบ \code{node\_modules} ให้เหลือกองเดียว
ที่ราก} แล้วทั้งสองฝั่ง \en{resolve} แพ็กเกจจากที่เดียวกัน

ด่าน 2 จึงหายไปเองโดยไม่ต้องทำอะไรเพิ่ม และด่าน 3 จะถูกบังคับให้แก้ เพราะ
\en{npm} จะบ่นทันทีถ้าสองฝั่งขอ \code{zod} คนละเวอร์ชันใหญ่

จากนั้นสร้าง \en{workspace} ตัวที่สามชื่อ \code{packages/contract} ไว้เก็บ
ไฟล์สัญญา --- ทั้งสองฝั่ง \code{import} จากชื่อแพ็กเกจ ไม่ใช่จากพาธสัมพัทธ์

\subsubsection{วิธีทำ ทีละขั้น}

\begin{enumerate}
\en{merge} \code{back-end-branch} เข้า \code{main} (0 \en{conflict}
      --- \code{front-end-branch} เข้าไปแล้ว)
\item ย้าย \code{backend-preview/backend/} $\to$ \code{backend/} ให้ชื่อโฟลเดอร์
      สม่ำเสมอ แล้วลบ \en{worktree} ทิ้งด้วย \code{git worktree remove}
\item สร้าง \code{package.json} ที่รากรีโป
\item สร้าง \code{packages/contract/} พร้อม \code{package.json} ของตัวเอง
\item ชี้ \code{autoSchemaFile} ของ \en{backend} ให้เขียนลงใน \en{package} นั้น
\item เพิ่ม \code{"@ulms/contract": "*"} ใน \code{dependencies} ของทั้งสองฝั่ง
\item \textbf{ลบ \code{node\_modules} และ \code{package-lock.json} ของทั้งสอง
      โฟลเดอร์ย่อยทิ้ง} แล้ว \code{npm install} ที่รากครั้งเดียว
\item ยก \code{zod} เป็น 4 ทั้งคู่ (ว-02) --- ขั้นนี้จะง่ายขึ้นเพราะเห็นชัดใน
      ที่เดียว
\end{enumerate}

\begin{lstlisting}[language={}]
{
  "name": "ulms",
  "private": true,
  "workspaces": ["frontend", "backend", "packages/*"],
  "scripts": {
    "dev:fe":    "npm run dev -w frontend",
    "dev:be":    "npm run start:dev -w backend",
    "typecheck": "npm run typecheck -w frontend && npm run typecheck -w backend"
  }
}
\end{lstlisting}

\begin{lstlisting}[language=bash]
# one install for the whole repo, from the root
rm -rf frontend/node_modules backend/node_modules
rm -f  frontend/package-lock.json backend/package-lock.json
npm install
\end{lstlisting}

\subsubsection{ปัญหาที่จะเกิด}

\begin{table}[H]
\centering\small
\begin{tabular}{L{4.4cm}L{5.0cm}L{5.0cm}}
\toprule
\sansall\bfseries ปัญหา & \sansall\bfseries อาการ & \sansall\bfseries กันไว้อย่างไร \\
\midrule
พาธใน \en{Prisma} พังตอนย้ายโฟลเดอร์ &
\code{schema.prisma} ตั้ง \code{output = "../src/generated/prisma"} แบบสัมพัทธ์
และ \code{prisma.config.ts} ชี้ \code{"prisma/schema.prisma"} &
รัน \en{Prisma} จากใน \en{workspace} เสมอ (\code{npm run -w backend ...})
ไม่ใช่จากราก และตรวจ \code{npx prisma generate} ให้ผ่านทันทีหลังย้าย \\
\addlinespace
\textbf{\en{phantom dependency}} &
\en{hoisting} ทำให้ \code{import} แพ็กเกจที่ไม่ได้ประกาศไว้ใน
\code{package.json} ของตัวเองได้ แล้วพังตอนแยกไป \en{deploy} &
ทุกครั้งที่ \code{import} ของใหม่ ให้ \code{npm i -w <ws> <pkg>} เสมอ
ห้ามพึ่งว่า ``มันก็เห็นอยู่แล้ว'' \\
\addlinespace
\en{lockfile} ใหญ่ขึ้นและ \en{conflict} บ่อยขึ้น &
\code{package-lock.json} ไฟล์เดียวสำหรับทุก \en{workspace} &
เวลา \en{conflict} \textbf{อย่าแก้ด้วยมือ} ให้เอาของฝั่งใดฝั่งหนึ่งแล้ว
\code{npm install} ใหม่เพื่อสร้างใหม่ \\
\addlinespace
\en{Vite} กับ \en{package} ที่เป็น \en{TypeScript} ดิบ &
\code{@ulms/contract} เป็นซอร์ส \code{.ts} ไม่ได้ \en{build} เป็น \code{.d.ts}
\en{Vite} อาจ \en{pre-bundle} ผิด &
ให้ \code{contract} \en{export} เป็นซอร์ส \code{.ts} ตรง ๆ และตั้ง
\code{optimizeDeps.exclude} ถ้าเจออาการแปลก \\
\addlinespace
ทุกคนต้องทำพร้อมกัน &
คนที่ยังไม่ลบ \code{node\_modules} เก่าจะเจออาการที่อธิบายไม่ได้ &
ประกาศล่วงหน้า ให้ทุกคน \code{git pull} แล้วรันชุดคำสั่งเดียวกัน \\
\bottomrule
\end{tabular}
\caption{ปัญหาของทางเลือก ก. --- สังเกตว่าทั้งหมดเป็นปัญหา ``ครั้งเดียวตอนย้าย'' ไม่ใช่ปัญหาที่กวนใจไปตลอด}
\end{table}

\subsubsection{สิ่งที่ควรรู้ก่อนลงมือ}

\begin{itemize}
\item \textbf{ไม่ต้องลงเครื่องมือใหม่} --- \en{npm} 10.9 ที่ทีมใช้อยู่รองรับ
      \en{workspaces} มาตั้งแต่ \en{npm} 7 และ \code{package-lock.json}
      ทั้งสองไฟล์เป็น \code{lockfileVersion 3} อยู่แล้ว
\item \textbf{\en{pnpm} ดีกว่าในเชิงเทคนิค แต่ไม่แนะนำตอนนี้} --- มันเข้มงวดกว่า
      และกัน \en{phantom dependency} ได้จริง แต่ทุกคนต้องลง \en{pnpm} เพิ่ม
      และเปลี่ยนคำสั่งที่ใช้ทุกวัน --- เป็นการเปลี่ยนสองเรื่องพร้อมกัน
      ซึ่งเวลาพังจะแยกไม่ออกว่าพังเพราะอะไร
\item \textbf{ทำเป็น \en{PR} เดียว คนเดียวทำ} คนอื่นหยุด \en{push} ระหว่างนั้น
      ประเมินเวลา 2--3 ชั่วโมง
\item \textbf{ต้องมีแผนถอย} --- ถ้าเกินเวลาที่ตั้งไว้แล้วยังไม่จบ ให้
      \code{git revert} แล้วไปทางเลือก ข. ก่อน อย่าปล่อยให้รีโปค้างอยู่ครึ่งทาง
      ข้ามคืน
\item \textbf{อย่าทำในสัปดาห์ที่มีเดดไลน์ส่งงาน}
\end{itemize}

\subsection{ทางเลือก ข. --- คัดลอกไฟล์ที่ถูกสร้าง}

\subsubsection{แนวคิด}

ปล่อยให้ \code{node\_modules} แยกกันเหมือนเดิม แล้วแก้ด่าน 2 ด้วยการ
\textbf{ทำให้ฝั่ง \en{frontend} มีของครบพอที่จะ \en{resolve} ไฟล์สัญญาได้เอง}
คือลง \code{zod} กับ \code{@trpc/server} เพิ่มในฝั่งนั้นด้วย

แล้วคัดลอกไฟล์สัญญาข้ามมาด้วยสคริปต์

\subsubsection{วิธีทำ ทีละขั้น}

\begin{enumerate}
\item \en{merge} สองสาขาเข้า \code{main} --- ทางเลือกนี้ \emph{ทำงานได้} โดยไม่
      \en{merge} เพราะ \en{worktree} ทำให้โค้ดอยู่บนดิสก์พร้อมกันแล้ว
      \textbf{แต่ \en{CI} จะตรวจไม่ได้ถ้าโค้ดยังอยู่คนละสาขา} จึงควร \en{merge} อยู่ดี
\item ตั้ง \code{autoSchemaFile} ใน \code{app.module.ts} ฝั่ง \en{backend}
\item \textbf{ลงแพ็กเกจเพิ่มฝั่ง \en{frontend} สามตัว} (ดูกล่องเตือนข้างล่าง)
\item เขียนสคริปต์ \code{scripts/sync-contract.sh} พร้อมโหมด \code{--check}
\item เพิ่ม \code{npm run sync:contract} และ \code{check:contract} ใน
      \code{package.json}
\item \textbf{ตั้ง \en{CI} หรือกติกาก่อน \en{merge}} ให้รัน
      \code{check:contract} กับ \code{tsc --noEmit} ทั้งสองฝั่ง
\end{enumerate}

\begin{warnbox}[title={\en{frontend} ต้องลงเพิ่มสามตัว ไม่ใช่ตัวเดียว --- จุดที่พลาดกันบ่อย}]
เพราะไฟล์สัญญาที่คัดลอกมามี \code{import} ของตัวเองอยู่ข้างใน ฝั่งที่รับต้องมี
ของครบด้วย

\begin{lstlisting}[language=bash]
cd frontend
npm i zod@^4               # must match the backend major exactly - see agenda 02
npm i @trpc/client         # to actually make calls
npm i -D @trpc/server      # types only — the contract file imports it
\end{lstlisting}

\code{@trpc/server} ฝั่ง \en{frontend} เป็น \code{devDependency} เพราะใช้แค่
\en{type} ตอน \en{compile} ไม่ได้ใช้ตอนรัน --- และเวอร์ชันต้องตรงกับฝั่ง
\en{backend} (ตอนนี้ 11.18) ไม่งั้นจะเจออาการ ``\en{Type} X ใช้แทน \en{Type} X
ไม่ได้'' แบบเดียวกับที่อธิบายไว้ในด่าน 2
\end{warnbox}

\begin{badbox}[title={ต้องเป็น \code{import type} เท่านั้น ห้ามลืมคำว่า \code{type}}]
ไฟล์สัญญามี \code{initTRPC.create()} ซึ่งเป็นโค้ดที่รันได้จริง ไม่ใช่ \en{type} ล้วน

\begin{lstlisting}
import type { AppRouter } from '@/server-types/appRouter';  // correct
import { AppRouter } from '@/server-types/appRouter';       // WRONG
\end{lstlisting}

บรรทัดล่างจะทำให้ \en{Vite} ลาก \code{@trpc/server} ทั้งก้อนเข้าไปใน \en{bundle}
ที่ส่งไปเบราว์เซอร์ --- \en{build} ผ่านปกติ ไม่มี \en{error} แต่ขนาดไฟล์จะโตขึ้น
โดยไม่มีใครสังเกต

กันด้วยการเปิดกฎ \en{ESLint} \code{consistent-type-imports}
\end{badbox}

\subsubsection{ปัญหาที่จะเกิด}

\begin{table}[H]
\centering\small
\begin{tabular}{L{4.4cm}L{5.0cm}L{5.0cm}}
\toprule
\sansall\bfseries ปัญหา & \sansall\bfseries อาการ & \sansall\bfseries กันไว้อย่างไร \\
\midrule
\textbf{สัญญาเก่าได้เงียบ ๆ} (ปัญหาหลัก) &
มีคนแก้ \en{schema} ฝั่ง \en{backend} แล้วลืมรันสคริปต์ \en{frontend} ยังใช้
สัญญาเก่าโดย \code{tsc} ไม่ฟ้อง &
\code{check:contract} ต้องเป็นเงื่อนไขก่อน \en{merge} --- \textbf{ถ้าไม่มีข้อนี้
ทางเลือก ข. แย่กว่าไม่มีสัญญาเลย} \\
\addlinespace
เวอร์ชัน \code{zod}/\code{@trpc/server} หลุดจากกันภายหลัง &
มีคนอัป \code{zod} ข้างเดียว แล้ว \en{type} เพี้ยนแบบอ่านไม่ออก &
เพิ่มการเทียบเวอร์ชันเข้าไปในสคริปต์ \code{--check} ด้วย ไม่ใช่ตรวจแค่เนื้อไฟล์ \\
\addlinespace
ไฟล์ \en{generated} เข้า \en{git} ทุกครั้ง &
\en{PR} มี \en{diff} ยาวมากจนรีวิวโค้ดจริงไม่เจอ &
ใส่ \code{.gitattributes} ทำเครื่องหมาย \code{linguist-generated} เพื่อให้
\en{GitHub} ยุบ \en{diff} ให้ \\
\addlinespace
\en{conflict} ในไฟล์ที่ถูก \en{generate} &
สอง \en{PR} แก้สัญญาพร้อมกัน &
\textbf{ห้ามแก้ \en{conflict} ด้วยมือ} ให้ \en{regenerate} แล้ว \en{commit} ทับ \\
\addlinespace
มี \code{zod} สองชุดบนดิสก์ตลอดไป &
ต่อให้เวอร์ชันตรงกัน ก็ยังมีโอกาสเจอ \en{error} \en{type} ซ้อน &
ยอมรับความเสี่ยงนี้ หรือย้ายไปทางเลือก ก. ในภายหลัง \\
\bottomrule
\end{tabular}
\caption{ปัญหาของทางเลือก ข. --- ต่างจาก ก. ตรงที่เป็นปัญหา ``กวนใจไปเรื่อย ๆ'' ไม่ใช่ปัญหาครั้งเดียว}
\end{table}

\subsubsection{สิ่งที่ควรรู้ก่อนลงมือ}

\begin{itemize}
\item \textbf{ตอนนี้รีโปยังไม่มี \en{CI} เลย} --- ไม่มีโฟลเดอร์ \code{.github/}
      แปลว่ากติกา ``ต้องผ่าน \code{check:contract} ก่อน \en{merge}'' ยังไม่มี
      อะไรบังคับได้จริง \textbf{เลือกทางนี้ = รับปากว่าจะตั้ง \en{CI} ด้วย}
      หรือไม่ก็ยอมรับว่าเป็นการตรวจด้วยคนล้วน ๆ
\item \textbf{ทางเลือกนี้ไม่ได้ห้ามย้ายไป ก. ทีหลัง} --- ย้ายได้ และตอนย้ายจะ
      ง่ายกว่าตอนนี้ เพราะสัญญาถูกแยกเป็นโฟลเดอร์ของตัวเองอยู่แล้ว
\item \textbf{สคริปต์ต้องเขียนให้ ``ล้มดัง''} --- ถ้าหาไฟล์ต้นทางไม่เจอ ต้อง
      \en{exit} พร้อมบอกสาเหตุ ไม่ใช่คัดลอกไฟล์เปล่าแล้วผ่าน
\item มีตัวอย่างสคริปต์เขียนไว้แล้วที่ \code{docs/trpc-example/scripts/sync-contract.sh}
      (ดูภาคผนวก ค)
\end{itemize}

\subsection{เทียบสองทางเลือกแบบตัดสินใจได้}

\begin{table}[H]
\centering\small
\begin{tabular}{L{4.0cm}L{5.2cm}L{5.2cm}}
\toprule
\sansall\bfseries เกณฑ์ & \sansall\bfseries ก. \en{workspaces} &
\sansall\bfseries ข. คัดลอกไฟล์ \\
\midrule
แก้ด่าน 2 ได้จริงไหม & \textbf{ใช่ ที่ต้นเหตุ} --- เหลือ \code{node\_modules} กองเดียว &
เลี่ยงได้ --- แต่ยังมีสองกอง ต้องดูแลเวอร์ชันเอง \\
\addlinespace
ต้นทุนตอนเริ่ม & สูง --- 2--3 ชม. ทุกคนหยุด \en{push} & ต่ำ --- เริ่มได้ในครึ่งวัน \\
\addlinespace
ต้นทุนต่อเนื่อง & \textbf{แทบไม่มี} & ต้องรันสคริปต์ทุกครั้ง + ต้องมี \en{CI} ตรวจ \\
\addlinespace
ถ้าทีมวินัยหลุด & ยังปลอดภัย \en{type} ตรงกันเสมอ &
\textbf{อันตราย} --- สัญญาเก่าเงียบ ๆ แย่กว่าไม่มีสัญญา \\
\addlinespace
ความเสี่ยงตอนทำ & มี --- อาจติดพาธ \en{Prisma} หรือ \en{Vite} & น้อยมาก \\
\addlinespace
ต้องมี \en{CI} ก่อนไหม & ไม่ต้อง & \textbf{ควรมี} (ตอนนี้ยังไม่มี) \\
\addlinespace
ย้ายไปอีกทางทีหลัง & ย้ายกลับยาก & ย้ายไป ก. ได้ง่าย \\
\bottomrule
\end{tabular}
\caption{แถวที่ชี้ขาดคือแถวที่สี่ --- คำถามจริงคือ ``ทีมเราจะรักษาวินัยรันสคริปต์ได้ตลอดไหม ในสัปดาห์ที่ใกล้เดดไลน์''}
\end{table}

\subsection{ข้อเสนอ}

\textbf{เลือก ก. --- \en{npm workspaces}} โดยมีเงื่อนไขสามข้อ

\begin{enumerate}
\item \en{merge} \code{back-end-branch} เข้า \code{main} ก่อน เพราะไม่มีต้นทุน
      (0 \en{conflict}) และเป็นสิ่งที่ต้องมีก่อนอยู่แล้วสำหรับทางเลือกนี้
\item ทำเป็น \en{PR} เดียว คนเดียวทำ คนอื่นหยุด \en{push} ระหว่างนั้น ---
      ตั้งเพดานเวลาไว้ 3 ชั่วโมง
\item \textbf{ต้องมีแผนถอย} ถ้าเกินเพดาน ให้ \code{git revert} แล้วไปทางเลือก ข.
      ชั่วคราว อย่าปล่อยให้รีโปค้างครึ่งทางข้ามคืน
\end{enumerate}

\textbf{เหตุผลที่เปลี่ยนมาเสนอ ก. แทน ข.} --- ร่างแรกของเอกสารนี้เสนอ ข. ด้วย
เหตุผลว่า ``เริ่มได้เร็วกว่า'' แต่พอไล่รายละเอียดครบทั้งสองทางแล้ว น้ำหนักเปลี่ยน

\begin{table}[H]
\centering\small
\begin{tabular}{L{4.6cm}L{10.0cm}}
\toprule
\sansall\bfseries ข้อค้นพบ & \sansall\bfseries ผลต่อการตัดสินใจ \\
\midrule
\textbf{รีโปยังไม่มี \en{CI} เลย} &
ทางเลือก ข. ทั้งหมดฝากความปลอดภัยไว้กับ \code{check:contract} ที่รันอัตโนมัติ
--- ซึ่งยังไม่มีอะไรรันมันได้ ต้นทุนจริงของ ข. จึงรวมค่าตั้ง \en{CI} เข้าไปด้วย
ไม่ได้ ``เริ่มได้เร็วกว่า'' อย่างที่คิดตอนแรก \\
\addlinespace
\textbf{ข. ยังต้องแก้ปัญหาเวอร์ชันอยู่ดี} &
ต้องลง \code{zod@4} และ \code{@trpc/server} เพิ่มฝั่ง \en{frontend} ให้ตรง
เวอร์ชันเป๊ะ แล้วต้องคอยกันไม่ให้หลุดจากกันตลอดไป ---
ก. ทำให้ปัญหานี้ \emph{หายไป} ไม่ใช่ \emph{ต้องคอยดูแล} \\
\addlinespace
\textbf{ก. ให้ประโยชน์ที่ ข. ให้ไม่ได้} &
\en{frontend} เอา \code{zod schema} ตัวเดียวกับที่ \en{backend} ใช้ \en{validate}
ไปตรวจฟอร์มได้เลย ไม่ต้องเขียนซ้ำใน \code{frontend/src/schemas/} ---
ซึ่งการเขียนซ้ำก็คือการเปิดช่องให้ความจริงสองชุดเลื่อนออกจากกันอีกรอบ
อันเป็นปัญหาเดิมที่เรากำลังแก้อยู่ \\
\addlinespace
\textbf{ความเสี่ยงของ ก. เป็นแบบ ``ครั้งเดียว''} &
ติดพาธ \en{Prisma} หรือ \en{tsconfig} ของ \en{Nest} --- เจอแล้วแก้จบ
มีทางแก้เขียนไว้ให้แล้วใน \code{trpc-example/README.md} หัวข้อกับดัก
ส่วนความเสี่ยงของ ข. เป็นแบบ ``กวนใจไปตลอด'' และพึ่งวินัยของคน \\
\bottomrule
\end{tabular}
\caption{สี่ข้อที่ทำให้ข้อเสนอเปลี่ยนจาก ข. เป็น ก. --- ข้อแรกคือข้อที่ชี้ขาด}
\end{table}

\begin{warnbox}[title={ทางเลือก ข. ยังเป็นทางถอยที่ใช้ได้ อย่าลบทิ้งจากความจำ}]
หัวข้อ ข. ข้างบนยังอยู่ในเอกสารเพราะสองเหตุผล: เป็นแผนถอยถ้าการย้ายไป
\en{workspaces} ติดปัญหาเกินเพดานเวลา และเป็นทางที่เลือกได้ถ้าที่ประชุมประเมินว่า
ตอนนี้ไม่ใช่จังหวะที่จะรื้อโครงรีโป

\code{trpc-example/STRUCTURE.md} หัวข้อ 8 สรุปไว้แล้วว่าถ้าสลับไป ข.
ต้องเปลี่ยนสี่จุดไหน
\end{warnbox}

\textbf{ถ้าไม่ตกลงวันนี้:} \en{backend} จะเขียน \en{router} ไปเรื่อย ๆ โดยที่
\en{frontend} ยังต้องเดา \en{type} เอง แล้วเมื่อถึงวันที่เชื่อมจริงจะพบว่าต้องแก้
ทั้งสองฝั่งพร้อมกัน --- ซึ่งคือสถานการณ์ที่ \en{tRPC} ควรจะป้องกันตั้งแต่แรก

\begin{outbox}
\begin{itemize}
\item จะ \en{merge} สองสาขาเข้า \code{main} เมื่อไร ใครทำ ---
\item เลือกทางเลือก ก. หรือ ข. ---
\item ถ้าเลือก ข. --- ใครรับผิดชอบตั้ง \en{CI} และเสร็จวันไหน ---
\item ถ้าเลือก ก. --- ใครทำ วันไหน และเวลาที่ยอมให้ใช้ก่อนตัดสินใจถอย ---
\item ตำแหน่งไฟล์สัญญาฝั่ง \en{backend} (ค่าของ \code{autoSchemaFile}) ---
\item ตำแหน่งปลายทางฝั่ง \en{frontend} ---
\item ยืนยันว่า \en{frontend} จะลง \code{zod@4}, \code{@trpc/client},
      \code{@trpc/server} (\en{dev}) และใครทำ
\end{itemize}
\end{outbox}

\section{ว-02 · ใช้ \code{zod} เวอร์ชันเดียวกันทั้งสองฝั่ง}

\textbf{คำถามที่ต้องตอบ:} จะยก \en{frontend} ขึ้น \code{zod} 4 หรือจะลด
\en{backend} ลงมา 3 และใครทำ เสร็จเมื่อไร

\textbf{ทำไมอยู่ระดับ \pzero{}:} ตอนนี้ \en{backend} ใช้ \code{zod} 4.4 ส่วน
\en{frontend} ใช้ 3.23 ไฟล์สัญญาที่ \code{nestjs-trpc} สร้างมีโค้ด \code{zod}
อยู่ข้างใน เมื่อ \en{frontend} เอาไป \en{compile} ด้วยรุ่นใหญ่คนละรุ่นจะได้
\en{error} ที่อ่านไม่ออกว่าสาเหตุคือเรื่องเวอร์ชัน --- และจะเสียเวลาไล่หาสาเหตุ
กันทั้งทีม

\textbf{ข้อเสนอ:} ยก \en{frontend} ขึ้น \code{zod} 4 เพราะ \en{backend} เขียน
โค้ดไปน้อยกว่ามาก และ \code{zod} 4 คือทิศทางที่ไลบรารีไปต่อ

\begin{warnbox}[title={งานนี้ไม่ใช่แค่เปลี่ยนตัวเลขใน \code{package.json}}]
ต้องตรวจของที่พึ่งพากันด้วย: \code{frontend/src/schemas/loan-request.schema.ts}
ที่เขียนไว้แล้ว, แพ็กเกจ \code{@hookform/resolvers} ที่เชื่อม
\code{react-hook-form} กับ \code{zod} และทุกจุดที่อ่านผลของ \code{safeParse}

\textbf{ต้องทำเป็น \en{PR} แยกให้จบก่อน} ไม่รวมกับงาน \en{tRPC} เพราะถ้ารวมกัน
แล้วพัง จะแยกไม่ออกว่าพังเพราะเรื่องไหน
\end{warnbox}

\textbf{ถ้าไม่ตกลงวันนี้:} คนแรกที่ลองเชื่อมสองฝั่งจะติดกำแพงนี้ และจะเสียเวลา
ครึ่งวันไปกับ \en{error} ที่ไม่ได้บอกสาเหตุจริง

\begin{outbox}
\begin{itemize}
\item เวอร์ชัน \code{zod} ที่ทั้งทีมใช้ ---
\item ใครทำการอัปเกรด และ \en{PR} นี้จะ \en{merge} วันไหน ---
\item ยืนยันแล้วหรือยังว่า \code{npm run typecheck} ฝั่ง \en{frontend} ผ่านหลังอัป
\end{itemize}
\end{outbox}

\section{ว-03 · รูปร่างของ \code{ctx.user} และวิธีส่งข้อมูลยืนยันตัวตน}

\textbf{คำถามที่ต้องตอบ:} \en{object} ที่ทุก \en{procedure} จะได้รับมีฟิลด์อะไรบ้าง
และ \en{token} เดินทางมาถึง \en{backend} ได้อย่างไร

\textbf{ทำไมอยู่ระดับ \pzero{}:} ทุก \en{procedure} ที่ต้องล็อกอินพึ่ง
\code{ctx.user} ถ้าเขียนไป 20 ตัวแล้วพบว่าขาดฟิลด์หนึ่ง ต้องกลับไปแก้ทั้ง 20 ตัว
พร้อมกับ \en{service} ที่เรียกใช้

\textbf{ข้อเสนอเรื่องการส่ง \en{token}:} เรื่องนี้ตัดสินใจไปแล้วโดยปริยาย ---
\code{frontend/src/lib/api-client.ts} ตั้ง \code{credentials: "include"} ไว้ทุก
\en{request} แปลว่าใช้ \en{httpOnly cookie} สิ่งที่เหลือคือทำให้ฝั่ง \en{backend}
รองรับ ซึ่งตอนนี้ \code{main.ts} ยังไม่มี \code{enableCors} เลย

\textbf{ข้อเสนอเรื่องฟิลด์:} เริ่มจากชุดนี้ และเพิ่มได้เฉพาะเมื่อมีเหตุผลชัดเจน

\begin{table}[H]
\centering\small
\begin{tabular}{L{3.4cm}L{4.0cm}L{7.0cm}}
\toprule
\sansall\bfseries ฟิลด์ & \sansall\bfseries มาจาก & \sansall\bfseries ใครใช้ \\
\midrule
\code{accountKey} & \code{AccountInfo.AccountKey} & ทุก \en{query} ที่กรองตามเจ้าของ \\
\code{role} & \code{RoleInfo.RoleName} & \en{middleware} ตรวจสิทธิ์ทุกตัว \\
\code{facultyKey} & \code{FacultyInfo} & กฎการยืมที่ผูกกับสังกัด \\
\code{creditScore} & \code{AccountInfo.UserCredit} &
  ใช้หาระดับเครดิต (\code{CreditTier}) เพื่อคำนวณ \textbf{จำนวนวันยืมสูงสุด}
  และจำนวนครั้งที่ต่ออายุได้ \\
\bottomrule
\end{tabular}
\caption{ชุดฟิลด์ขั้นต่ำที่เสนอ --- ทุกฟิลด์ต้องตอบได้ว่ามีใครใช้ ไม่งั้นอย่าใส่}
\end{table}

\begin{keybox}[title={เครดิตต่ำ \emph{ไม่ได้} แปลว่ายืมไม่ได้ --- แปลว่ายืมได้สั้นลง}]
ข้อนี้ต้องเข้าใจตรงกันทั้งทีมก่อนออกแบบ \en{output} เพราะมันเปลี่ยนว่า
\code{loan.create} ควรตอบอะไรกลับมา

ใน \code{schema.prisma} ระดับเครดิตเชื่อมกับกฎการยืมผ่านตาราง
\code{BorrowConstraints} ซึ่งเก็บ \code{MaxBorrowDate} (จำนวนวันยืมสูงสุด) กับ
\code{MaxExtendTime} (จำนวนครั้งที่ต่ออายุได้) ต่อคู่ ``กฎการยืม $\times$
ระดับเครดิต'' --- \textbf{ไม่มีฟิลด์ไหนเลยที่บอกว่าเครดิตเท่าไรถึงห้ามยืม}

สิ่งที่กันไม่ให้ยืมคือคนละเรื่องกัน: ตาราง \code{Eligibility} (สิทธิ์ตามสังกัด
และบทบาท) และ \code{MinimumAuthorityLevel} --- ทั้งคู่ไม่เกี่ยวกับคะแนนเครดิต
\end{keybox}

\begin{warnbox}[title={\code{creditScore} ใน \code{ctx} เสี่ยงเป็นข้อมูลเก่า}]
ถ้า \en{cron} หักเครดิตรายวันระหว่างที่ผู้ใช้กำลังกรอกคำขออยู่ ค่าที่ค้างใน
\code{ctx} จะไม่ตรงกับฐานข้อมูล ผลคือหน้าจอบอกว่ายืมได้ 14 วัน แต่ระบบให้จริง 7 วัน

ต้องตกลงว่า \en{procedure} ที่ \emph{คำนวณจำนวนวันยืมจริง} (เช่น
\code{loan.create}) อ่านคะแนนสดจากฐานข้อมูลเสมอ ส่วนค่าใน \code{ctx}
ใช้แสดงผลอย่างเดียว
\end{warnbox}

\textbf{ถ้าไม่ตกลงวันนี้:} ไม่มีใครเขียน \en{procedure} ที่ต้องล็อกอินได้เลย
ซึ่งคือเกือบทั้งหมดของระบบ

\begin{outbox}
\begin{itemize}
\item รายการฟิลด์ใน \code{ctx.user} (สรุปให้ครบ ห้ามค้าง ``ไว้ค่อยเพิ่ม'') ---
\item ยืนยันใช้ \en{httpOnly cookie} ---
\item รายการ \en{origin} ที่จะอนุญาตใน \en{CORS} ทั้งตอนพัฒนาและตอนใช้จริง ---
\item อายุของ \en{session} และพฤติกรรมเมื่อหมดอายุ ---
\item ใครทำ \en{context} + \en{CORS} และเสร็จเมื่อไร
\end{itemize}
\end{outbox}

\section{ว-04 · ใครเป็นเจ้าของสัญญา และการเปลี่ยนสัญญาทำอย่างไร}

\textbf{คำถามที่ต้องตอบ:} เมื่อมีคนอยากเพิ่มฟิลด์หรือเปลี่ยนชื่อ \en{procedure}
กระบวนการคืออะไร และใครมีสิทธิ์ตัดสินขั้นสุดท้าย

\textbf{ทำไมอยู่ระดับ \pzero{}:} นี่เป็นวาระเดียวในภาคนี้ที่ไม่ใช่เรื่องเทคนิค
แต่เป็นวาระที่ทำให้อีกสามวาระมีผลจริง สัญญาที่ไม่มีเจ้าของจะถูกแก้โดยใครก็ได้
โดยไม่มีใครรู้ และภายในสองสัปดาห์จะไม่มีใครเชื่อถือมันอีก

\textbf{ข้อเสนอ:} กำหนดสามอย่าง

\begin{enumerate}
\item \textbf{เจ้าของ} --- หนึ่งคนต่อหนึ่งโดเมน (\code{item}, \code{loan},
      \code{approval}, \code{admin}) เจ้าของเป็นคนอนุมัติการเปลี่ยน \en{schema}
      ของโดเมนนั้น ไม่ใช่คนเดียวคุมทั้งหมด เพราะจะกลายเป็นคอขวด
\item \textbf{การแบ่งชนิดการเปลี่ยน} --- แยกให้ชัดว่าอันไหนแจ้งอย่างเดียวพอ
      อันไหนต้องคุยก่อน
\item \textbf{ช่องทางแจ้ง} --- ที่เดียว ทุกคนเห็น ไม่ใช่แชตส่วนตัว
\end{enumerate}

\begin{table}[H]
\centering\small
\begin{tabular}{L{4.2cm}L{4.6cm}L{5.6cm}}
\toprule
\sansall\bfseries ชนิดการเปลี่ยน & \sansall\bfseries ตัวอย่าง & \sansall\bfseries ต้องทำอะไร \\
\midrule
เพิ่มฟิลด์ที่ไม่บังคับ &
เพิ่ม \code{note?} ใน \en{output} &
แจ้งในช่องทางกลาง ไม่ต้องรออนุมัติ \\
\addlinespace
เพิ่ม \en{procedure} ใหม่ &
เพิ่ม \code{loan.extend} &
แจ้ง + เพิ่มแถวในตารางสัญญา \\
\addlinespace
\textbf{เปลี่ยนชื่อ / ลบฟิลด์ / เปลี่ยนชนิด} &
\code{dueAt} เป็น \code{dueDate} &
\textbf{ต้องคุยกับอีกฝั่งก่อน} และแก้พร้อมกันใน \en{PR} เดียว \\
\addlinespace
เปลี่ยนความหมายโดยรูปร่างเหมือนเดิม &
\code{status} เดิมนับ ``รอรับของ'' เป็น \code{pending} เปลี่ยนเป็นสถานะใหม่ &
\textbf{อันตรายที่สุด} --- \en{compiler} จับไม่ได้ ต้องคุยและต้องมีคนทดสอบ \\
\bottomrule
\end{tabular}
\caption{สี่ชนิดการเปลี่ยน --- แถวล่างสุดคือชนิดที่ระบบ \en{type} ช่วยอะไรไม่ได้เลย จึงต้องใช้กระบวนการของคนล้วน ๆ}
\end{table}

\textbf{ถ้าไม่ตกลงวันนี้:} จะมีคนแก้ \en{schema} แล้วอีกฝั่งพังโดยไม่รู้ตัว
เกิดขึ้นแน่นอน คำถามคือเกิดกี่ครั้งและเสียเวลาไปเท่าไร

\begin{outbox}
\begin{itemize}
\item เจ้าของแต่ละโดเมน ---
\item ช่องทางแจ้งการเปลี่ยนสัญญา ---
\item ยืนยันว่าการเปลี่ยนแบบ \en{breaking} ต้องแก้สองฝั่งใน \en{PR} เดียว ---
\item ใครดูแลตารางสัญญาให้ตรงกับโค้ด
\end{itemize}
\end{outbox}

\clearpage

% =============================================================================
\part{\pone{} ต้องจบภายในสัปดาห์นี้}
% =============================================================================

\textbf{หกวาระนี้ไม่บล็อกการเริ่มต้น แต่บล็อก ``ความสม่ำเสมอ'' --- เขียน
\en{procedure} ตัวแรกได้โดยไม่ต้องตกลง แต่ถ้าตกลงตอนมี 30 ตัวแล้วจะต้องแก้ทั้ง 30
เวลาที่เสนอ: 35 นาทีในรอบนี้ ที่เหลือแยกคุยเป็นรอบย่อย}

\section{ว-05 · โครงและการตั้งชื่อ \en{router} กับ \en{procedure}}

\textbf{คำถามที่ต้องตอบ:} จัดกลุ่ม \en{procedure} ตามอะไร และใช้คำกริยาชุดไหน

\textbf{ทำไมอยู่ระดับ \pone{}:} เปลี่ยนชื่อทีหลังต้องแก้ทุกจุดที่เรียกใช้ทั้งสองฝั่ง
และเป็นการเปลี่ยนที่ไม่ได้ประโยชน์อะไรกลับมานอกจากความสม่ำเสมอ --- งานที่ทีมมักไม่ยอมทำ
แล้วปล่อยให้ชื่อมั่วต่อไปจนจบโครงการ

\textbf{ข้อเสนอ 1 --- จัดกลุ่มตามโดเมน ไม่ใช่ตามบทบาทผู้ใช้}

\begin{table}[H]
\centering\small
\begin{tabular}{L{6.8cm}L{7.6cm}}
\toprule
\sansall\bfseries แบบที่เสนอ (ตามโดเมน) & \sansall\bfseries แบบที่ไม่เสนอ (ตามบทบาท) \\
\midrule
\code{approval.decide} ตัวเดียว ตรวจสิทธิ์ด้วย \en{middleware} &
\code{staff.approveLoan} กับ \code{supervisor.approveLoan} สองตัว \\
\addlinespace
ตรรกะการอนุมัติอยู่ที่เดียว แก้ครั้งเดียว &
ตรรกะเดียวกันถูกคัดลอกสองที่ แก้ที่หนึ่งลืมอีกที่ \\
\addlinespace
เพิ่มบทบาทใหม่ = เพิ่ม \en{middleware} &
เพิ่มบทบาทใหม่ = เพิ่ม \en{router} ทั้งชุด \\
\bottomrule
\end{tabular}
\caption{จัดกลุ่มตามโดเมนชนะเพราะบทบาทเปลี่ยนบ่อยกว่าโดเมน --- ในระบบนี้มี 4--5 บทบาทที่ยังอาจปรับได้อีก}
\end{table}

โดเมนที่เสนอ: \code{auth}, \code{item}, \code{loan}, \code{reservation},
\code{approval}, \code{appeal}, \code{credit}, \code{inspection},
\code{notification}, \code{report}, \code{admin}

\textbf{ข้อเสนอ 2 --- ชุดคำกริยาที่ใช้ซ้ำทั้งระบบ}

\code{list} · \code{getById} · \code{create} · \code{update} · \code{cancel} ·
\code{decide} · \code{confirm} --- ห้ามปน \code{getAll} กับ \code{fetchList}
กับ \code{findMany} ในความหมายเดียวกัน

\begin{outbox}
\begin{itemize}
\item รายชื่อโดเมน (\en{router}) ที่ตกลง ---
\item ชุดคำกริยามาตรฐาน ---
\item ยืนยันว่าตรวจสิทธิ์ด้วย \en{middleware} ไม่ใช่แยก \en{router} ตามบทบาท
\end{itemize}
\end{outbox}

\section{ว-06 · รูปแบบ \en{error} และรายการรหัสข้อผิดพลาดทางธุรกิจ}

\textbf{คำถามที่ต้องตอบ:} เมื่อ ``ของถูกคนอื่นจองตัดหน้าไปแล้ว'' \en{backend}
ส่งอะไรกลับมา และ \en{frontend} แปลงเป็นข้อความที่ไหน

\textbf{ทำไมอยู่ระดับ \pone{}:} \en{frontend} ต้องเขียน \en{UI} รองรับแต่ละกรณี
รายการรหัสจึงเป็น \textbf{ส่วนหนึ่งของ \en{output}} ไม่ใช่รายละเอียดปลีกย่อย ---
ถ้าไม่มีรายการตั้งแต่ต้น \en{frontend} จะเขียนได้แค่ ``เกิดข้อผิดพลาด'' แล้วต้อง
กลับมาแก้ทีหลังทุกหน้า

\textbf{ข้อเสนอ:} สองชั้น --- ชั้นนอกใช้รหัสมาตรฐานของ \en{tRPC}
(\code{UNAUTHORIZED}, \code{FORBIDDEN}, \code{NOT\_FOUND}, \code{CONFLICT},
\code{BAD\_REQUEST}) ชั้นในใส่รหัสธุรกิจของเราเอง

\begin{table}[H]
\centering\small
\begin{tabular}{L{4.4cm}L{2.6cm}L{7.4cm}}
\toprule
\sansall\bfseries รหัสธุรกิจ & \sansall\bfseries รหัส \en{tRPC} &
\sansall\bfseries \en{frontend} ต้องทำอะไร \\
\midrule
\code{NOT\_ELIGIBLE} & \code{FORBIDDEN} & บอกว่าสังกัด/บทบาทนี้ยืมของชิ้นนี้ไม่ได้ \\
\code{ITEM\_UNAVAILABLE} & \code{CONFLICT} & แสดงวันที่พร้อมให้ยืมถัดไป \\
\code{SLOT\_TAKEN} & \code{CONFLICT} & \en{refetch} สล็อตแล้วให้เลือกใหม่ \\
\code{LOAN\_PERIOD\_EXCEEDS\_LIMIT} & \code{BAD\_REQUEST} &
  บอกเพดานวันยืมของระดับเครดิตปัจจุบัน แล้วให้เลือกวันใหม่ \\
\code{EXTENSION\_QUOTA\_EXCEEDED} & \code{FORBIDDEN} & บอกจำนวนครั้งที่ใช้ไปแล้วจากเพดานของระดับเครดิต \\
\code{ALREADY\_DECIDED} & \code{CONFLICT} & \en{refetch} คิว --- มีคนตัดสินไปก่อนแล้ว \\
\code{APPEAL\_WINDOW\_CLOSED} & \code{FORBIDDEN} & บอกวันที่หมดเขต \\
\bottomrule
\end{tabular}
\caption{ร่างรายการรหัส --- ยังไม่ครบ ต้องเติมให้ครบตอนแปลงเอกสารรายการเรียกใช้งานเป็นตารางสัญญา สังเกตว่า \emph{ไม่มี} รหัสว่า ``เครดิตต่ำจนยืมไม่ได้'' เพราะกฎแบบนั้นไม่มีในระบบ}
\end{table}

\begin{warnbox}[title={คำถามที่ยังไม่มีใครตอบ: เครดิตต่ำแล้วขอยืมนานเกินเพดาน จะ ``ปฏิเสธ'' หรือ ``ให้เท่าที่ได้''}]
เพราะเครดิตกำหนดแค่ \emph{ความยาว} ของการยืม ไม่ใช่สิทธิ์ยืม จึงเกิดทางแยกที่
กระทบทั้ง \en{error} และ \en{output} ของ \code{loan.create}

\begin{itemize}
\item \textbf{ก. ปฏิเสธ} --- โยน \code{LOAN\_PERIOD\_EXCEEDS\_LIMIT} พร้อมเพดาน
      แล้วให้ผู้ใช้เลือกวันใหม่ · ชัดเจน ไม่มีเซอร์ไพรส์ แต่เพิ่มขั้นตอนให้ผู้ใช้
\item \textbf{ข. ตัดให้อัตโนมัติ} --- อนุมัติแต่ย่นวันคืนให้เท่าเพดาน ·
      ลื่นไหลกว่า แต่ \textbf{\en{output} ต้องบอกกลับมาให้ชัดว่าได้กี่วันจริง
      และถูกย่นเพราะอะไร} ไม่งั้นผู้ใช้จะคืนของสาย
\end{itemize}

ข้อเสนอ: เลือก \textbf{ก.} สำหรับ \en{MVP} เพราะแสดงผลง่ายกว่าและไม่ต้องออกแบบ
ช่องทางแจ้งเตือนเพิ่ม --- แต่ไม่ว่าจะเลือกทางไหน \en{output} ของ
\code{loan.create} ต้องมีวันครบกำหนดจริงเสมอ ห้ามให้ \en{frontend} คำนวณเอง
\end{warnbox}

\begin{badbox}[title={ห้าม \en{backend} ส่งข้อความภาษาไทยมาให้แสดงตรง ๆ}]
โปรเจกต์ลง \code{i18next} และ \code{react-i18next} ไว้แล้ว แปลว่าตั้งใจรองรับ
หลายภาษา ถ้า \en{backend} ส่ง ``ขอยืมได้ไม่เกิน 7 วัน'' มาแล้ว \en{frontend}
เอาไปแสดงเลย ระบบหลายภาษาจะใช้ไม่ได้ทันที และข้อความจะกระจายอยู่ในโค้ดสองที่

\en{backend} ส่ง \textbf{รหัส} กับ \textbf{ข้อมูลประกอบ} (เช่น เพดานวันยืมของ
ระดับเครดิตปัจจุบัน จำนวนวันที่ขอมา) แล้วให้ \en{frontend} ประกอบเป็นประโยค
\end{badbox}

\begin{outbox}
\begin{itemize}
\item ยืนยันโครงสร้าง \en{error} สองชั้น ---
\item รายการรหัสธุรกิจฉบับเริ่มต้น และใครเป็นคนเติมให้ครบ ---
\item ยืนยันว่าข้อความภาษาไทยอยู่ฝั่ง \en{frontend} เท่านั้น ---
\item เลือก ก. ปฏิเสธ หรือ ข. ตัดวันยืมให้อัตโนมัติ เมื่อขอเกินเพดานของระดับเครดิต ---
\item รูปแบบข้อมูลประกอบที่แนบไปกับ \en{error}
\end{itemize}
\end{outbox}

\section{ว-07 · รูปแบบการแบ่งหน้าและการกรอง}

\textbf{คำถามที่ต้องตอบ:} \en{procedure} ที่คืนรายการ จะรับพารามิเตอร์ชุดไหน
และคืนอะไรกลับมา

\textbf{ทำไมอยู่ระดับ \pone{}:} เอกสารรายการเรียกใช้งานมีรายการที่ต้องแบ่งหน้า
อย่างน้อย 5 จุด ถ้าแต่ละคนออกแบบเอง \en{frontend} จะต้องเขียน \en{component}
ตารางคนละแบบสำหรับแต่ละหน้า

\begin{table}[H]
\centering\small
\begin{tabular}{L{2.6cm}L{5.4cm}L{6.4cm}}
\toprule
\sansall\bfseries แบบ & \sansall\bfseries หน้าตา & \sansall\bfseries เหมาะกับ \\
\midrule
\textbf{ตามเลขหน้า} &
รับ \code{page, pageSize, sort, q} คืน \code{items, total, page, pageSize} &
\textbf{หน้าจอฝั่ง \en{staff}/\en{admin} ที่ต้องเห็นว่ามีทั้งหมดกี่หน้า และกระโดดไปหน้าที่ต้องการได้} \\
\addlinespace
ตามตัวชี้ (\en{cursor}) &
รับ \code{cursor, limit} คืน \code{items, nextCursor} &
เลื่อนลงเรื่อย ๆ บนมือถือ ประสิทธิภาพดีกว่าเมื่อข้อมูลเยอะมาก \\
\bottomrule
\end{tabular}
\caption{สองแบบ --- ข้อเสนอคือแบบเลขหน้า เพราะหน้าจอส่วนใหญ่ในระบบนี้เป็นตารางงานของเจ้าหน้าที่}
\end{table}

\textbf{ข้อเสนอ:} ใช้แบบเลขหน้าทั้งระบบ กำหนด \code{pageSize} ปริยาย 20
และเพดาน 100 (กันคนยิงขอทีเดียวหมื่นแถว)

\begin{outbox}
\begin{itemize}
\item แบบที่เลือก และชื่อฟิลด์ที่ตกลง ---
\item ค่าปริยายและเพดานของ \code{pageSize} ---
\item รูปแบบการส่งเงื่อนไขกรอง (แยกเป็นฟิลด์ หรือรวมเป็น \en{object} \code{filters})
\end{itemize}
\end{outbox}

\section{ว-08 · รูปแบบของวันที่และเวลา}

\textbf{คำถามที่ต้องตอบ:} ส่ง \code{Date} จริงผ่าน \code{superjson} หรือส่งเป็น
สตริง \en{ISO}

\textbf{ทำไมอยู่ระดับ \pone{}:} \code{schema.prisma} มีฟิลด์ \code{DateTime}
กระจายอยู่ 9 ตารางหลัก และระบบนี้ ``เวลา'' คือแก่นของตรรกะธุรกิจ --- กำหนดคืน
การเลยกำหนด การหักเครดิตรายวัน อายุบทลงโทษ ถ้าตกลงไม่ตรงกัน บั๊กจะไปโผล่ที่
การคำนวณค่าปรับ

\textbf{ข้อเสนอ:} สตริง \en{ISO} เพราะดีบักง่าย (เปิด \en{Network tab} แล้วอ่านออก
ทันที) ไม่ต้องตั้งค่าให้ตรงกันสองฝั่ง และทีมมี \code{date-fns} อยู่แล้ว

\begin{warnbox}[title={สิ่งที่สำคัญกว่าการเลือกแบบไหน คือห้ามมีข้อยกเว้น}]
\en{procedure} ที่ส่ง \code{Date} จริงปนกับ \en{procedure} ที่ส่งสตริง คือ
สถานการณ์ที่แย่ที่สุดในสามแบบ เพราะ \en{frontend} จะไม่รู้ว่าต้อง \code{parseISO}
เมื่อไร --- และ \en{type} จะบอกว่าถูกทั้งสองแบบ
\end{warnbox}

นอกจากรูปแบบแล้วต้องตกลงอีกสองเรื่องที่ลืมกันบ่อย: \textbf{เขตเวลา}
(เก็บเป็น \en{UTC} แล้วแปลงตอนแสดงผลใช่ไหม) และ \textbf{ค่าที่เป็นวันล้วน}
เช่นวันครบกำหนดคืน --- เป็นวันที่อย่างเดียวหรือมีเวลาด้วย และถ้ามีเวลา
กำหนดคืน ``วันที่ 20'' หมายถึงกี่โมง

\begin{outbox}
\begin{itemize}
\item รูปแบบที่เลือก ---
\item เขตเวลาที่เก็บและที่แสดง ---
\item ฟิลด์ที่เป็นวันล้วน จัดการอย่างไร ---
\item เวลาตัดของ ``ครบกำหนดคืน''
\end{itemize}
\end{outbox}

\section{ว-09 · กฎการแปลงชื่อฟิลด์ และห้ามส่ง \en{Prisma model} ออกไปตรง ๆ}

\textbf{คำถามที่ต้องตอบ:} \en{output} จะพูดภาษาฐานข้อมูลหรือภาษา \en{frontend}

\textbf{ทำไมอยู่ระดับ \pone{}:} \code{schema.prisma} ตั้งชื่อแบบ
\code{PascalCase} (\code{UserFName}, \code{AccountKey}) ส่วน
\code{frontend/src/types/domain.ts} ใช้ \code{camelCase} (\code{firstName},
\code{id}) ถ้าไม่ตกลงตอนนี้ \en{procedure} ที่เขียนก่อนจะส่งชื่อแบบหนึ่ง
ตัวที่เขียนทีหลังส่งอีกแบบ

\textbf{ข้อเสนอ:} \en{output} เป็นภาษา \en{frontend} เสมอ (\code{camelCase})
และห้าม \code{return} \en{Prisma model} ตรง ๆ ในทุกกรณี

\begin{badbox}[title={เหตุผลที่ห้ามไม่ใช่แค่เรื่องความสวยงามของชื่อ}]
ตาราง \code{AccountInfo} มีฟิลด์ \code{HashedPassword} อยู่ในตารางเดียวกับชื่อ
และอีเมล การ \code{findUnique} แล้วส่งกลับตรง ๆ คือการส่ง \en{hash} รหัสผ่าน
ออกไปที่เบราว์เซอร์ --- และไม่มีอะไรในระบบเตือนเลย

ประโยชน์อีกข้อคือฐานข้อมูลเปลี่ยนได้โดยหน้าเว็บไม่พัง ซึ่งสำคัญกับรีโปนี้เป็นพิเศษ
เพราะ \en{schema} ยังอยู่ในช่วงที่อาจปรับอีก
\end{badbox}

\begin{outbox}
\begin{itemize}
\item ยืนยันว่า \en{output} ใช้ \code{camelCase} ---
\item ยืนยันกฎ ``ห้ามส่ง \en{Prisma model} ตรง ๆ'' และใครตรวจตอนรีวิว ---
\item ที่เก็บ \en{output schema} กลางที่นำไป \code{.pick()}/\code{.extend()} ต่อได้
\end{itemize}
\end{outbox}

\section{ว-10 · สถานะส่งออกเป็นสตริงคงที่ ไม่ใช่เลข \en{key}}

\textbf{คำถามที่ต้องตอบ:} \en{frontend} จะได้รับสถานะเป็นอะไร

\textbf{ทำไมอยู่ระดับ \pone{}:} \code{schema.prisma} \textbf{ไม่มี \en{enum}
เลยสักตัว} --- สถานะทุกชนิดเป็นตารางแยก (\code{ApproveStatus},
\code{ResourceStatus}, \code{ConditionType}, \code{NotificationType},
\code{PenaltyType} และแม้แต่ \code{RoleInfo}) ขณะที่ \en{frontend} นิยามไว้
เป็นสตริง เช่น \code{"available" | "reserved" | "borrowed"}

\begin{badbox}[title={ส่งเลข \en{primary key} ของสถานะออกไป = บั๊กที่จะโผล่ตอนย้ายเครื่อง}]
\code{ApproveStatusKey = 2} หมายถึง ``อนุมัติแล้ว'' บนฐานข้อมูลเครื่องหนึ่ง
แต่ถ้าอีกเครื่อง \code{seed} คนละลำดับ เลข 2 อาจกลายเป็น ``ปฏิเสธ''

โค้ด \code{if (status === 2)} จะคอมไพล์ผ่านทุกครั้งและผิดเงียบ ๆ ---
เป็นบั๊กที่จะเจอตอนสาธิตงาน ไม่ใช่ตอนพัฒนา
\end{badbox}

\textbf{ข้อเสนอ:} \en{output} ใช้ \code{z.enum([...])} ที่เป็นสตริงคงที่เสมอ
และต้องตกลง \textbf{รายการค่าที่เป็นไปได้ของแต่ละสถานะ} ให้ครบในวาระนี้ เพราะ
\en{frontend} ต้องเขียน \en{UI} ครอบคลุมทุกค่า

\begin{outbox}
\begin{itemize}
\item ยืนยันว่าสถานะส่งเป็นสตริง ---
\item รายการค่าของ \code{ApproveStatus} ---
\item รายการค่าของ \code{ResourceStatus} ---
\item รายการค่าของ \code{RoleInfo} ---
\item ใครดูแลให้ตารางในฐานข้อมูลกับ \code{z.enum} ไม่หลุดจากกัน
\end{itemize}
\end{outbox}

\clearpage

% =============================================================================
\part{\ptwo{} ตกลงตอนเริ่มทำ \en{feature} นั้น}
% =============================================================================

\textbf{ห้าวาระนี้ \emph{ไม่ต้องคุยในการประชุมรอบนี้} เอกสารระบุไว้เพื่อให้ทุกคน
เห็นว่ามันไม่ได้ถูกลืม และรู้ว่าจะถูกหยิบมาคุยเมื่อไร --- การเขียนสิ่งที่ยัง
ไม่ตัดสินใจลงกระดาษ คือสิ่งที่ทำให้ ``ผลัดไปก่อน'' ต่างจาก ``ลืม''}

\section{ว-11 · ขั้นตอนการอัปโหลดรูป}

\textbf{คำถาม:} รูปตอนรับของและตอนคืนของเดินทางอย่างไร เก็บที่ไหน

\textbf{ทำไมผลัดได้:} \en{tRPC} ส่งไฟล์ไม่ได้อยู่แล้ว จึงต้องใช้เส้นทาง
\en{REST} แยก ซึ่ง \code{api-client.ts} มี \code{uploadFile} รองรับไว้แล้ว
--- โครงสร้างส่วนนี้ไม่ขึ้นกับการตัดสินใจเรื่องสัญญา

\textbf{คุยเมื่อ:} เริ่มทำ \en{feature} รับ--คืนของ

\textbf{ประเด็นที่ต้องเตรียมไว้:} ที่เก็บไฟล์ · การขอ \en{URL} ชั่วคราว ·
เพดานขนาดและชนิดไฟล์ · การจัดการไฟล์ที่อัปแล้วแต่ไม่มีแถวใน \code{Images} ชี้ถึง
(ผู้ใช้ปิดเบราว์เซอร์กลางคัน) · การบีบอัดรูปก่อนอัปจากมือถือ

\section{ว-12 · ความถี่ \en{polling} จริง และ \code{staleTime}}

\textbf{คำถาม:} ตัวเลขในเอกสารรายการเรียกใช้งานใช้ได้จริงไหม

\textbf{ทำไมผลัดได้:} เปลี่ยนตัวเลขเป็นการแก้บรรทัดเดียวฝั่ง \en{frontend}
ไม่กระทบสัญญา

\textbf{คุยเมื่อ:} หลังต่อหน้าจอที่ \en{polling} หน้าแรกได้แล้ว และวัดภาระจริงได้

\textbf{ประเด็น:} 10--15 วินาทีสำหรับของคงเหลือเหมาะกับจำนวนผู้ใช้จริงหรือไม่ ·
\code{staleTime} ยาวสำหรับข้อมูลหลัก (หมวดหมู่ \en{tier} ระดับความเสียหาย) ·
หยุด \en{polling} เมื่อผู้ใช้สลับแท็บออกไป

\section{ว-13 · การรวมคำขอ เวลาหมดอายุ และการลองใหม่}

\textbf{คำถาม:} ใช้ \code{httpBatchLink} ทั้งหมด หรือแยกเส้นสำหรับ \en{query} ที่ถี่

\textbf{ทำไมผลัดได้:} เปลี่ยน \en{link} เป็นการแก้ไฟล์ตั้งค่าไฟล์เดียว

\textbf{คุยเมื่อ:} พบว่าหน้าใดหน้าหนึ่งช้าผิดปกติ

\textbf{ประเด็น:} \en{query} ที่ \en{poll} ทุก 10 วินาทีไปรวมก้อนกับรายงานสถิติ
จะถูกลากให้ช้าตาม · จะลองใหม่กี่ครั้งเมื่อเครือข่ายหลุด · \en{mutation}
ต้องไม่ลองใหม่อัตโนมัติ (ยืมซ้ำสองครั้ง)

\section{ว-14 · ข้อความ \en{error} อยู่ฝั่งไหน}

\textbf{คำถาม:} แผนผังจากรหัสธุรกิจไปเป็นประโยคเก็บที่ไหน

\textbf{ทำไมผลัดได้:} เมื่อ ว-06 ตกลงแล้วว่าส่งเป็นรหัส ที่เหลือเป็นงานภายในของ
\en{frontend} ล้วน ๆ

\textbf{คุยเมื่อ:} เริ่มทำหน้าจอที่มี \en{error} หลายกรณี

\textbf{ประเด็น:} รวมเข้ากับ \code{lib/error-messages.ts} ที่มีอยู่แล้ว หรือย้าย
เข้าไฟล์แปลของ \code{i18next} · ข้อความปริยายเมื่อเจอรหัสที่ยังไม่ได้แปล

\section{ว-15 · การบันทึกและตามรอยปัญหา}

\textbf{คำถาม:} เมื่อผู้ใช้แจ้งว่า ``กดยืมแล้วขึ้น \en{error}'' จะหาสาเหตุอย่างไร

\textbf{ทำไมผลัดได้:} เพิ่มทีหลังได้โดยไม่กระทบสัญญา

\textbf{คุยเมื่อ:} ใกล้ส่งงาน หรือเมื่อเริ่มมีผู้ใช้จริงทดลอง

\textbf{ประเด็น:} \en{log} ชื่อ \en{procedure} เวลาที่ใช้ และรหัสข้อผิดพลาด ·
รหัสอ้างอิงต่อ \en{request} ที่แสดงให้ผู้ใช้บอกกลับมาได้ · ระวังอย่า \en{log}
ข้อมูลส่วนบุคคล

\clearpage

% =============================================================================
\part{\pthree{} ผลัดไปหลัง \en{MVP} ได้}
% =============================================================================

\textbf{สี่วาระนี้สำคัญในระบบที่ใช้งานจริง แต่ไม่กระทบการส่งงาน และที่สำคัญคือ
\emph{แก้ทีหลังได้โดยไม่ต้องรื้อของเดิม} --- เขียนไว้เพื่อไม่ให้หายไปจากความทรงจำของทีม}

\section{ว-16 · การอัปเดตแบบเรียลไทม์}

ทีมเลือก \en{polling} ไปแล้วในเอกสารรายการเรียกใช้งาน ซึ่งเป็นการตัดสินใจที่
เหมาะกับขนาดของระบบและเวลาที่มี \en{tRPC} รองรับ \en{subscription} ได้ถ้าอนาคต
ต้องการ และการย้ายจาก \en{polling} ไป \en{subscription} ไม่กระทบรูปร่างของ
\en{output} ที่ตกลงกันไว้ --- นี่คือเหตุผลที่ผลัดได้อย่างปลอดภัย

\section{ว-17 · การจำกัดอัตราการเรียกและการป้องกันการยิงถี่}

\en{query} ที่ \en{poll} ทุก 10 วินาทีคูณจำนวนผู้ใช้พร้อมกันเป็นภาระที่คาดเดาได้
ส่วนที่คาดเดาไม่ได้คือคนที่ยิงตรงเข้ามา เพิ่มการจำกัดอัตราทีหลังได้ที่ชั้น
\en{middleware} โดยไม่แตะสัญญา

\section{ว-18 · การทำเวอร์ชันของสัญญาอย่างเป็นทางการ}

ตราบใดที่ยังส่งพร้อมกันทั้งสองฝั่ง ไม่ต้องมีเวอร์ชัน เรื่องนี้จำเป็นเมื่อมี
\en{client} ที่อัปเดตไม่พร้อมกัน เช่นแอปมือถือที่ผู้ใช้ยังไม่กดอัป ---
ยังไม่ใช่สถานการณ์ของโครงการนี้

\section{ว-19 · การทดสอบสัญญาอัตโนมัติ}

\code{tsc --noEmit} ที่ตกลงใน ว-01 ครอบคลุมความเสี่ยงหลักไปแล้ว ส่วนการเขียน
เทสต์ที่ยิง \en{procedure} จริงแล้วตรวจ \en{output} กับ \en{schema} เป็นชั้นถัดไป
ที่เพิ่มได้เมื่อมีเวลา --- \en{backend} มี \code{jest} กับ \code{supertest}
ติดตั้งไว้แล้ว

\clearpage

% =============================================================================
\part{วิธีดำเนินการประชุม}
% =============================================================================

\textbf{คำถามที่ภาคนี้ตอบ: ทำอย่างไรให้ 90 นาทีนี้จบด้วยมติ ไม่ใช่จบด้วย
``เดี๋ยวไปคุยกันต่อ''}

\section{ตารางเวลาที่เสนอ}

\begin{table}[H]
\centering\small
\begin{tabular}{L{1.6cm}L{7.4cm}L{5.4cm}}
\toprule
\sansall\bfseries เวลา & \sansall\bfseries วาระ & \sansall\bfseries หมายเหตุ \\
\midrule
5 นาที & ทบทวนแผนที่ในรูปที่ \ref{fig:dep} & ให้ทุกคนเห็นตรงกันว่าจะคุยอะไรบ้าง \\
\addlinespace
25 นาที & \pzero{} ว-01 ขนส่งสัญญา & วาระที่ยาวที่สุด --- เผื่อเวลาสำหรับคำถาม ``\en{merge} ก็จบแล้วไม่ใช่เหรอ'' ไว้แล้ว \\
5 นาที & \pzero{} ว-02 เวอร์ชัน \code{zod} & น่าจะจบเร็ว มีข้อเสนอชัดอยู่แล้ว \\
20 นาที & \pzero{} ว-03 \code{ctx.user} + \en{cookie} & ต้องได้รายการฟิลด์ที่ครบ \\
10 นาที & \pzero{} ว-04 เจ้าของสัญญา & เรื่องคน ไม่ใช่เรื่องเทคนิค \\
\addlinespace
20 นาที & \pone{} ว-05 ถึง ว-10 รวบคุย & ถ้าเวลาไม่พอ ให้ยกไปรอบย่อย --- อย่ายืมเวลาจาก \pzero{} \\
\addlinespace
5 นาที & ทวนมติ + ผู้รับผิดชอบ + กำหนดเสร็จ & \textbf{ห้ามข้าม} \\
\bottomrule
\end{tabular}
\caption{90 นาที --- สังเกตว่า \pzero{} กินเวลา 55 นาที จากทั้งหมด 90 นาที ซึ่งเป็นสัดส่วนที่ตั้งใจ}
\end{table}

\section{กติกาที่ทำให้ประชุมจบเป็นมติ}

\begin{enumerate}
\item \textbf{ทุกวาระต้องได้ครบสามอย่าง} --- ข้อตกลง · ผู้รับผิดชอบ · วันเสร็จ
      ขาดข้อใดข้อหนึ่งแปลว่ายังไม่จบวาระ
\item \textbf{มีคนจับเวลา} และมีสิทธิ์ตัดบทได้จริง
\item \textbf{ถ้าเถียงกันเกิน 5 นาทีโดยไม่มีข้อมูลใหม่} ให้เลือกข้อเสนอในเอกสาร
      ไปก่อน แล้วบันทึกว่าเป็นการตัดสินใจแบบทบทวนได้ --- การตัดสินใจที่กลับได้
      ไม่ควรใช้เวลาเท่าการตัดสินใจที่กลับไม่ได้
\item \textbf{เรื่องที่ไม่อยู่ในภาค 1--2 ห้ามคุย} จดลงรายการค้างแล้วไปต่อ
\item \textbf{ผู้เข้าร่วมต้องอ่านภาค 1 มาก่อน} ---
      ถ้ายังไม่มีใครอ่าน ให้เลื่อนประชุมดีกว่าใช้เวลาครึ่งหนึ่งไปกับการอธิบาย
\end{enumerate}

\begin{keybox}[title={สิ่งที่วัดว่าประชุมนี้สำเร็จ}]
กรอกกล่องฟ้า ``ผลลัพธ์ที่ต้องจดลงบันทึก'' ของวาระ ว-01 ถึง ว-04 ได้ครบทุกช่อง
--- ไม่ใช่ ``คุยครบทุกวาระ''
\end{keybox}

\clearpage

% =============================================================================
\appendix
\part*{ภาคผนวก}
\addcontentsline{toc}{part}{ภาคผนวก}
% =============================================================================

\makeatletter
\newcommand{\thaialph}[1]{\ifcase\value{#1}\or ก\or ข\or ค\or ง\or จ\or ฉ\fi}
\makeatother
\renewcommand{\thesection}{\thaialph{section}}

\section{แบบฟอร์มบันทึกมติ}

พิมพ์หน้านี้ออกมาหรือเปิดคู่ไว้ระหว่างประชุม กรอกให้ครบทุกช่องก่อนปิดวาระ

\begin{longtable}{L{1.3cm}L{4.8cm}L{4.0cm}L{2.0cm}L{2.0cm}}
\toprule
\sansall\bfseries วาระ & \sansall\bfseries ข้อตกลง &
\sansall\bfseries รายละเอียด/เงื่อนไข & \sansall\bfseries ผู้รับผิดชอบ &
\sansall\bfseries วันเสร็จ \\
\midrule
\endfirsthead
\toprule
\sansall\bfseries วาระ & \sansall\bfseries ข้อตกลง &
\sansall\bfseries รายละเอียด/เงื่อนไข & \sansall\bfseries ผู้รับผิดชอบ &
\sansall\bfseries วันเสร็จ \\
\midrule
\endhead
\bottomrule
\endfoot
ว-01 & \vspace{9mm} & & & \\
ว-02 & \vspace{9mm} & & & \\
ว-03 & \vspace{9mm} & & & \\
ว-04 & \vspace{9mm} & & & \\
ว-05 & \vspace{9mm} & & & \\
ว-06 & \vspace{9mm} & & & \\
ว-07 & \vspace{9mm} & & & \\
ว-08 & \vspace{9mm} & & & \\
ว-09 & \vspace{9mm} & & & \\
ว-10 & \vspace{9mm} & & & \\
\end{longtable}

\section{เช็กลิสต์ก่อนเข้าประชุม}

\begin{table}[H]
\centering\small
\begin{tabular}{L{0.6cm}L{8.6cm}L{5.2cm}}
\toprule
\sansall\bfseries \# & \sansall\bfseries สิ่งที่ต้องเตรียม & \sansall\bfseries ใคร \\
\midrule
1 & อ่านภาค 1 ของเอกสารนี้ (\pzero{} สี่วาระ) & ทุกคน \\
2 & อ่านภาค 4 ของ \code{trpc-guide.tex} (ความเสี่ยง 8 ข้อ) & อย่างน้อยฝ่ายเทคนิค \\
3 & เปิดดู \code{frontend/package.json} กับ \code{backend-preview/backend/package.json} เทียบเวอร์ชัน \code{zod} จริง & คนเสนอ ว-02 \\
4 & เตรียมรายการฟิลด์ที่คิดว่า \code{ctx.user} ต้องมี & \en{backend} \\
5 & เตรียมรายการหน้าจอที่ต้องใช้ข้อมูลผู้ใช้ & \en{frontend} \\
6 & เปิด \code{schema.prisma} ค้างไว้ที่ตารางสถานะ (\code{ApproveStatus}, \code{ResourceStatus}) & คนนำ ว-10 \\
7 & เปิด \code{docs/trpc-example/} ค้างไว้ --- ตัวอย่างโค้ดของทุกวาระ (ภาคผนวก ค) & ทุกคน \\
8 & เปิด \code{trpc-example/STRUCTURE.md} เตรียมกรอกช่องผู้รับผิดชอบ & ผู้จดบันทึก \\
9 & เตรียมกระดาษ/ไฟล์สำหรับกรอกแบบฟอร์มภาคผนวก ก & ผู้จดบันทึก \\
\bottomrule
\end{tabular}
\caption{เช็กลิสต์ --- ข้อ 1 เป็นข้อที่ตัดสินว่าประชุมจะได้ผลหรือไม่}
\end{table}

\section{โค้ดตัวอย่างของผลลัพธ์ --- \code{docs/trpc-example/}}

มีโฟลเดอร์ \code{docs/trpc-example/} ที่เขียน \textbf{หน้าตาของโค้ดเมื่อประชุม
จบแล้ว} ไว้ให้ครบทุกวาระใน \pzero{} และ \pone{} --- เป็นไฟล์จริงที่ก๊อปวางไปใช้ได้
ไม่ใช่ \en{pseudocode}

จุดประสงค์คือให้ดูปลายทางก่อนถกเถียง เพราะการเถียงว่า ``แบ่งหน้าแบบไหนดี''
บนกระดาษกินเวลากว่าการดูโค้ดสองแบบวางข้างกันมาก

\begin{table}[H]
\centering\footnotesize
\begin{tabular}{L{1.1cm}L{6.6cm}L{6.9cm}}
\toprule
\sansall\bfseries วาระ & \sansall\bfseries ไฟล์ที่แสดงผลลัพธ์ &
\sansall\bfseries ถ้าที่ประชุมเลือกอย่างอื่น \\
\midrule
ว-01 & \code{package.json} (\en{workspace root})\newline\code{packages/contract/}\newline\code{backend/src/app.module.ts} &
เลือก ข. $\to$ ลบ \code{packages/contract/} ย้าย \en{schema} กลับเข้า
\code{backend/} แล้วเพิ่ม \code{scripts/sync-contract.sh} \\
\addlinespace
ว-02 & \code{packages/contract/package.json} ประกาศ \code{zod} ที่เดียว &
เลือกลด \en{backend} เป็น \code{zod} 3 $\to$ แก้ \code{z.email()} กับ \code{z.iso.datetime()} \\
\addlinespace
ว-03 & \code{backend/src/trpc/context.ts}\newline\code{backend/src/trpc/auth.middleware.ts}\newline\code{backend/src/main.ts} &
แก้ฟิลด์ใน \code{TrpcUser} ที่เดียว แล้ว \en{compiler} จะฟ้องทุกจุดที่ต้องตามแก้ \\
\addlinespace
ว-04 & \code{process/CONTRACT-OWNERS.md}\newline\code{process/pull\_request\_template.md} &
กรอกชื่อเจ้าของในตารางระหว่างประชุม \\
\addlinespace
ว-05 & \code{backend/src/loan/loan.router.ts}\newline\code{backend/src/auth/auth.router.ts} &
เปลี่ยนชุดคำกริยา $\to$ แก้ชื่อเมท็อดในทุก \en{router} \\
\addlinespace
ว-06 & \code{packages/contract/src/errors/error-codes.ts}\newline\code{common/errors/business-error.ts}\newline\code{frontend/src/lib/error-messages.ts} &
เลือก ``ตัดวันยืมอัตโนมัติ'' แทนปฏิเสธ $\to$ แก้บล็อกใน \code{loan.service.ts}
ตามคอมเมนต์ที่เขียนไว้ให้แล้ว \\
\addlinespace
ว-07 & \code{packages/contract/src/schemas/pagination.schema.ts} &
เลือก \en{cursor} $\to$ เปลี่ยน \code{paginated()} เป็นคืน \code{nextCursor} \\
\addlinespace
ว-08 & \code{packages/contract/src/schemas/datetime.schema.ts} &
เลือก \code{superjson} $\to$ ลบไฟล์นี้ แล้วตั้ง \en{transformer} สองฝั่ง \\
\addlinespace
ว-09 & \code{packages/contract/src/schemas/user.schema.ts}\newline\code{common/mappers/user.mapper.ts} &
--- (แทบไม่มีทางเลือกอื่นที่ปลอดภัย) \\
\addlinespace
ว-10 & \code{packages/contract/src/schemas/status.schema.ts} &
เติมรายการค่าให้ตรงกับตารางสถานะจริงในฐานข้อมูล \\
\addlinespace
รวม & \code{CONTRACT.md} ตารางสัญญาที่คนอ่านได้\newline\code{frontend/.../use-loans.ts} ตัวอย่างการใช้จริง &
อัปเดตทุกครั้งที่เพิ่ม \en{procedure} \\
\addlinespace
โครงสร้าง & \code{STRUCTURE.md} --- ของจริงควรอยู่ที่ไหน ใครรับผิดชอบ &
มีโครงเป้าหมาย \textbf{ทั้งสองแบบ} (ก. และ ข.) วางเทียบกันไว้แล้ว
พร้อมตารางเจ้าของที่เว้นให้กรอกในห้อง \\
\bottomrule
\end{tabular}
\caption{แผนที่จากวาระไปหาไฟล์ --- ทุกไฟล์มีคอมเมนต์หัวไฟล์บอกว่าตอบวาระไหน สมมติว่าเลือกอะไร และปลายทางจริงอยู่พาธไหน}
\end{table}

\begin{warnbox}[title={โค้ดชุดนั้นเขียนบน \emph{ข้อเสนอ} ยังไม่ใช่มติ}]
ทุกไฟล์เขียนบนสมมติฐานว่าที่ประชุมเลือกตามข้อเสนอในเอกสารฉบับนี้ ถ้าที่ประชุม
ตัดสินต่างออกไป คอลัมน์ขวาสุดของตารางบอกไว้แล้วว่าต้องแก้ไฟล์ไหน

โค้ดชุดนั้นเขียนบน \textbf{ทางเลือก ก.} ตามข้อเสนอข้างต้น ถ้าที่ประชุมเลือก ข.
ดู \code{STRUCTURE.md} หัวข้อ 8

สถานะการตรวจสอบ: ไฟล์ \code{.ts} ทั้ง 21 ไฟล์ผ่าน \code{tsc --strict} โดยไม่มี
\en{syntax error} เหลือ แต่ \textbf{ยังไม่ได้พิสูจน์ว่า \en{type} ตรงกันจริง}
เพราะเครื่องที่สร้างยังไม่ได้ลง \code{node\_modules} ของทั้งสองโปรเจกต์ ---
ต้องรัน \code{tsc --noEmit} ซ้ำหลังคัดลอกไปวางในโปรเจกต์จริง
\end{warnbox}

\clearpage

\section{สรุปทั้ง 19 วาระในหน้าเดียว}

\begin{longtable}{L{1.2cm}L{1.3cm}L{7.0cm}L{4.6cm}}
\toprule
\sansall\bfseries วาระ & \sansall\bfseries ระดับ & \sansall\bfseries เรื่อง &
\sansall\bfseries ตกลงเมื่อไร \\
\midrule
\endfirsthead
\toprule
\sansall\bfseries วาระ & \sansall\bfseries ระดับ & \sansall\bfseries เรื่อง &
\sansall\bfseries ตกลงเมื่อไร \\
\midrule
\endhead
\bottomrule
\endfoot
ว-01 & \pzero{} & ไฟล์สัญญาเดินทางข้ามฝั่งอย่างไร & ก่อนออกจากห้อง \\
ว-02 & \pzero{} & เวอร์ชัน \code{zod} ตรงกันทั้งสองฝั่ง & ก่อนออกจากห้อง \\
ว-03 & \pzero{} & รูปร่าง \code{ctx.user} + \en{cookie}/\en{CORS} & ก่อนออกจากห้อง \\
ว-04 & \pzero{} & เจ้าของสัญญาและวิธีขอเปลี่ยน & ก่อนออกจากห้อง \\
\addlinespace
ว-05 & \pone{} & ชื่อ \en{router} และ \en{procedure} & ในสัปดาห์นี้ \\
ว-06 & \pone{} & รูปแบบ \en{error} + รหัสธุรกิจ & ในสัปดาห์นี้ \\
ว-07 & \pone{} & การแบ่งหน้าและการกรอง & ในสัปดาห์นี้ \\
ว-08 & \pone{} & รูปแบบวันที่และเขตเวลา & ในสัปดาห์นี้ \\
ว-09 & \pone{} & แปลงชื่อฟิลด์ + ห้ามส่ง \en{Prisma model} & ในสัปดาห์นี้ \\
ว-10 & \pone{} & สถานะเป็นสตริงคงที่ & ในสัปดาห์นี้ \\
\addlinespace
ว-11 & \ptwo{} & ขั้นตอนอัปโหลดรูป & เริ่มทำ \en{feature} รับ--คืน \\
ว-12 & \ptwo{} & ความถี่ \en{polling} จริง & หลังต่อหน้าแรกที่ \en{poll} \\
ว-13 & \ptwo{} & \en{batching} / \en{timeout} / \en{retry} & เมื่อเจอหน้าที่ช้า \\
ว-14 & \ptwo{} & ข้อความ \en{error}  & หน้าจอที่มี \en{error} หลายกรณี \\
ว-15 & \ptwo{} & การ \en{log} และตามรอยปัญหา & ใกล้ส่งงาน \\
\addlinespace
ว-16 & \pthree{} & \en{realtime} / \en{subscription} & หลัง \en{MVP} \\
ว-17 & \pthree{} & จำกัดอัตราการเรียก & หลัง \en{MVP} \\
ว-18 & \pthree{} & เวอร์ชันของสัญญา & เมื่อมี \en{client} ที่อัปไม่พร้อมกัน \\
ว-19 & \pthree{} & เทสต์สัญญาอัตโนมัติ & เมื่อมีเวลา \\
\end{longtable}

\vfill
\begin{center}
\begin{tikzpicture}
  \draw[rule, line width=0.6pt] (0,0) -- (\textwidth,0);
\end{tikzpicture}
\end{center}
{\sansall\footnotesize\color{muted}
เอกสารนี้เป็นวาระประชุม ไม่ใช่บันทึกการประชุม --- เมื่อประชุมเสร็จ
ให้บันทึกมติลงแบบฟอร์มภาคผนวก ก แล้วเก็บไว้คู่กัน \\
เหตุผลเบื้องหลังของทุกวาระ และที่มาของข้อเท็จจริงที่อ้างถึง
อยู่ในเอกสารคู่กัน \code{docs/trpc-guide.tex} ภาคผนวก ง \par}

\end{document}
